\documentclass[11pt]{report}

\usepackage[spanish]{babel}
\usepackage[utf8]{inputenc}
\usepackage[T1]{fontenc}
\usepackage{amsmath, amssymb, amsthm}
\usepackage{geometry}
\geometry{margin=2.5cm}
\usepackage{amsfonts}
\usepackage{enumitem}
\usepackage{xcolor}
\usepackage{hyperref}

% --- Configuración de entornos ---
\newtheorem{ex}{Ejercicio}
\newenvironment{solution}{\begin{proof}[Solución]}{\end{proof}}

\newtheorem{ejercicio}{Ejercicio}[chapter]
\newenvironment{solucion}{\color{blue}\textbf{Solución:}}{}


% --- Comando para productos internos ---
\newcommand{\inner}[2]{( #1, #2 )_V}

% --- Comando para normas ---
\newcommand{\norma}[1]{\Vert #1 \Vert_V}

% --- Comando para el espacio L² ---
\newcommand{\Ldos}{L^2}

% --- Comando para espacios de Sobolev ---
\newcommand{\Huno}{H^1}
\newcommand{\HunoCero}{H^1_0}

% --- Comando para la diferencial de Fréchet ---
\newcommand{\diff}{\mathrm{d}}

%\newcommand{\R}{\mathbb{R}}

% ---------- Macros ----------
\newcommand{\R}{\mathbb{R}}
\DeclareMathOperator*{\supx}{sup}
\newcommand{\dom}{\operatorname{dom}}
\newcommand{\epi}{\operatorname{epi}}
\newcommand{\cl}{\operatorname{cl}}
\newcommand{\ri}{\operatorname{ri}}
\newcommand{\indic}{I}
\newcommand{\supp}{\sigma}
\newcommand{\subdiff}{\partial}
\newcommand{\conv}{\operatorname{conv}}


\title{Prácticos de análisis convexo -- versi\'on 1.00}
\author{Franz Chouly \and Federico Correa}
\date{\today}


\begin{document}

\maketitle
\tableofcontents

%%%%%%%%%%%%%%%%%%%%%%%%%%%%%%%%%%%%%%%%%%%%%%%%%%

\chapter{Repaso}

\section{Espacios de Hilbert}

\noindent
Primero, unos ejercicios generales sobre espacios de Hilbert. En lo siguiente, $V$ es un espacio de Hilbert sobre $\mathbb{R}$.

\begin{ex}
Sea $V$ un espacio de Hilbert. Demostrar que para todo $u, v \in V$, se cumple la desigualdad de Cauchy-Schwarz:
\begin{equation*}
|\inner{u}{v}| \leq \norma{u} \norma{v}.
\end{equation*}
\end{ex}

%\begin{solution}
%Para cualquier $t \in \mathbb{R}$, consideremos el vector $u + t v \in V$. Entonces:
%\begin{equation*}
%0 \leq \norma{u + t v}^2 = \norma{u}^2 + 2 t \inner{u}{v} + t^2 \norma{v}^2.
%\end{equation*}
%Esto define un polinomio cuadrático en $t$ que siempre es no negativo. Por lo tanto, su discriminante debe ser menor o igual a cero:
%\begin{equation*}
%(2 \inner{u}{v})^2 - 4 \norma{u}^2 \norma{v}^2 \leq 0.
%\end{equation*}
%Simplificando, obtenemos:
%\begin{equation*}
%|\inner{u}{v}|^2 \leq \norma{u}^2 \norma{v}^2,
%\end{equation*}
%y tomando raíces cuadradas, se sigue la desigualdad de Cauchy-Schwarz.
%\end{solution}

\begin{ex}
Demostrar que en un espacio de Hilbert $V$, dos vectores $u, v \in V$ son ortogonales si y solo si:
\begin{equation*}
\norma{u + v}^2 = \norma{u}^2 + \norma{v}^2.
\end{equation*}
\end{ex}

%\begin{solution}
%($\Rightarrow$) Si $u$ y $v$ son ortogonales, entonces $\inner{u}{v} = 0$. Por lo tanto:
% \begin{equation*}
% \norma{u + v}^2 = \norma{u}^2 + 2 \inner{u}{v} + \norma{v}^2 = \norma{u}^2 + \norma{v}^2.
% \end{equation*}
%
%($\Leftarrow$) Supongamos que $\norma{u + v}^2 = \norma{u}^2 + \norma{v}^2$. Desarrollando el cuadrado de la norma:
% \begin{equation*}
% \norma{u}^2 + 2 \inner{u}{v} + \norma{v}^2 = \norma{u}^2 + \norma{v}^2,
% \end{equation*}
% lo que implica $2 \inner{u}{v} = 0$, es decir, $\inner{u}{v} = 0$. Por lo tanto, $u$ y $v$ son ortogonales.
%\end{solution}

\begin{ex}
Sea $V$ un espacio de Hilbert y sea $W \subset V$ un subespacio vectorial. Demostrar que $W$ es un subespacio de Hilbert si y solo si $W$ es cerrado en $V$.
\end{ex}

%\begin{solution}
%($\Rightarrow$) Si $W$ es un subespacio de Hilbert, entonces $W$ es completo con la norma inducida por $V$. Como $W$ es completo y $V$ es un espacio métrico, $W$ es cerrado en $V$.

%($\Leftarrow$) Supongamos que $W$ es cerrado en $V$. Como $W$ es un subespacio vectorial de $V$, la restricción del producto interno de $V$ a $W$ define un producto interno en $W$.

%Además, como $W$ es cerrado en $V$, toda sucesión de Cauchy en $W$ converge a un elemento de $W$ (ya que $V$ es completo). Por lo tanto, $W$ es completo y, en consecuencia, un subespacio de Hilbert.

%\end{solution}

\begin{ex}
Sea $V$ un espacio de Hilbert y sea $(e_n)_{n \in \mathbb{N}}$ una sucesión ortonormal en $V$. Demostrar que:
\begin{equation*}
\sum_{n=1}^\infty \inner{u}{e_n}^2 \leq \norma{u}^2 \quad \forall u \in V.
\end{equation*}
\end{ex}

%\begin{solution}
%Para cualquier $N \in \mathbb{N}$, definimos el vector:
%\begin{equation*}
%u_N = \sum_{n=1}^N \inner{u}{e_n} e_n.
%\end{equation*}
%Como $(e_n)$ es ortonormal, se tiene:
%\begin{equation*}
%\norma{u_N}^2 = \sum_{n=1}^N |\inner{u}{e_n}|^2.
%\end{equation*}
%Además, $u - u_N$ es ortogonal a $u_N$, por lo que:
%\begin{equation*}
%\norma{u}^2 = \norma{u_N}^2 + \norma{u - u_N}^2 \geq \norma{u_N}^2 = \sum_{n=1}^N |\inner{u}{e_n}|^2.
%\end{equation*}
%Tomando límite cuando $N \to \infty$, obtenemos:
%\begin{equation*}
%\sum_{n=1}^\infty |\inner{u}{e_n}|^2 \leq \norma{u}^2.
%\end{equation*}
%\end{solution}

\begin{ex}
Sea $V$ un espacio de Hilbert. Demostrar que toda sucesión ortonormal $(e_n)_{n \in \mathbb{N}}$ en $V$ es una base ortonormal si y solo si el subespacio generado por $(e_n)$ es denso en $V$.
\end{ex}

%\begin{solution}

%($\Rightarrow$) Si $(e_n)$ es una base ortonormal, entonces para todo $u \in V$, se tiene:
% \begin{equation*}
%u = \sum_{n=1}^\infty \inner{u}{e_n} e_n.
% \end{equation*}
% Esto implica que el subespacio generado por $(e_n)$ es denso en $V$ (de hecho, igual a $V$).

%($\Leftarrow$) Supongamos que el subespacio generado por $(e_n)$ es denso en $V$. Para todo $u \in V$ y $\epsilon > 0$, existe una combinación lineal finita $v = \sum_{n=1}^N \alpha_n e_n$ tal que $\norma{u - v} < \epsilon$.

%Como $(e_n)$ es ortonormal, los coeficientes $\alpha_n$ pueden elegirse como $\alpha_n = \inner{u}{e_n}$ para minimizar $\norma{u - v}$. Por lo tanto, $(e_n)$ es una base ortonormal.

%\end{solution}

\begin{ex}
Sea $V$ un espacio de Hilbert. Demostrar que si $(u_n)$ es una sucesión en $V$ que converge débilmente a $u \in V$, entonces:
\begin{equation*}
\norma{u} \leq \liminf_{n \to \infty} \norma{u_n}.
\end{equation*}
\end{ex}

%\begin{solution}
%Por la convergencia débil, para todo $v \in V$, se tiene $\inner{u_n}{v} \to \inner{u}{v}$. En particular, para $v = u$, se tiene:
%\begin{equation*}
%\inner{u_n}{u} \to \inner{u}{u} = \norma{u}^2.
%\end{equation*}
%Por otra parte, por la desigualdad de Cauchy-Schwarz:
%\begin{equation*}
%|\inner{u_n}{u}| \leq \norma{u_n} \norma{u}.
%\end{equation*}
%Tomando límite inferior cuando $n \to \infty$, obtenemos:
%\begin{equation*}
%\norma{u}^2 \leq \liminf_{n \to \infty} \norma{u_n} \norma{u}.
%\end{equation*}
%Si $u \neq 0$, podemos dividir ambos lados por $\norma{u}$ para obtener:
%\begin{equation*}
%\norma{u} \leq \liminf_{n \to \infty} \norma{u_n}.
%\end{equation*}
%Si $u = 0$, la desigualdad es trivial.
%\end{solution}

\begin{ex}
Sea $V$ un espacio de Hilbert. Demostrar que si $(u_n)$ es una sucesión en $V$ que converge fuertemente a $u \in V$, entonces $(u_n)$ también converge débilmente a $u$.
\end{ex}

%\begin{solution}
%Si $(u_n)$ converge fuertemente a $u$, entonces $\norma{u_n - u} \to 0$. Para cualquier $v \in V$, por la desigualdad de Cauchy-Schwarz:
%\begin{equation*}
%|\inner{u_n}{v} - \inner{u}{v}| = |\inner{u_n - u}{v}| \leq \norma{u_n - u} \norma{v} \to 0.
%\end{equation*}
%Por lo tanto, $\inner{u_n}{v} \to \inner{u}{v}$ para todo $v \in V$, lo que significa que $(u_n)$ converge débilmente a $u$.
%\end{solution}

%\begin{ex}
%Sea $V$ un espacio de Hilbert. Demostrar que el producto interno $\inner{\cdot}{\cdot}: V \times V \to \mathbb{R}$ es una función continua.
%\end{ex}

%\begin{solution}
%Para demostrar que $\inner{\cdot}{\cdot}$ es continua, debemos mostrar que para cualquier sucesión $(u_n, v_n) \to (u, v)$ en $V \times V$, se tiene $\inner{u_n}{v_n} \to \inner{u}{v}$. 

%Escribimos:
%\begin{align*}
%|\inner{u_n}{v_n} - \inner{u}{v}| &= |\inner{u_n}{v_n} - \inner{u_n}{v} + \inner{u_n}{v} - \inner{u}{v}| \
%&\leq |\inner{u_n}{v_n - v}| + |\inner{u_n - u}{v}| \
%&\leq \norma{u_n} \norma{v_n - v} + \norma{u_n - u} \norma{v}.
%\end{align*}
%Como $(u_n)$ y $(v_n)$ convergen fuertemente a $u$ y $v$ respectivamente, se tiene que $\norma{u_n}$ está acotada (porque $(u_n)$ converge) y $\norma{v_n - v} \to 0$, $\norma{u_n - u} \to 0$. Por lo tanto, $\inner{u_n}{v_n} \to \inner{u}{v}$.
%\end{solution}

\begin{ex}
Sea $V$ un espacio de Hilbert y sea $u \in V$ fijo. Demostrar que la función $f: V \to \mathbb{R}$ definida por $f(v) = \inner{u}{v}$ es un funcional lineal continuo.
\end{ex}

%\begin{solution}
%
%
%
%
%
%Linealidad: Para todo $v_1, v_2 \in V$ y $\alpha, \beta \in \mathbb{R}$, se tiene:
% \begin{align*}
% f(\alpha v_1 + \beta v_2) &= \inner{u}{\alpha v_1 + \beta v_2} \
% &= \alpha \inner{u}{v_1} + \beta \inner{u}{v_2} \
% &= \alpha f(v_1) + \beta f(v_2).
% \end{align*}
% Por lo tanto, $f$ es lineal.



%Continuidad: Por la desigualdad de Cauchy-Schwarz:
% \begin{equation*}
% |f(v)| = |\inner{u}{v}| \leq \norma{u} \norma{v}.
% \end{equation*}
% Por lo tanto, $f$ es continua (de hecho, $f$ es Lipschitz con constante $\norma{u}$).
%\end{solution}



%\begin{ex}
%Sea $V$ un espacio de Hilbert. Demostrar que toda base ortonormal $(e_i){i \in I}$ de $V$ satisface la igualdad de Parseval:
%\begin{equation*}
%\norma{u}^2 = \sum{i \in I} |\inner{u}{e_i}|^2 \quad \forall u \in V.
%\end{equation*}
%\end{ex}

%\begin{solution}
%Sea $u \in V$. Por definición de base ortonormal, se tiene:
%\begin{equation*}
%u = \sum_{i \in I} \inner{u}{e_i} e_i.
%\end{equation*}
%Tomando el producto interno con $u$ a ambos lados:
%\begin{equation*}
%\inner{u}{u} = \sum_{i \in I} |\inner{u}{e_i}|^2 + \inner{u - \sum_{i \in I} \inner{u}{e_i} e_i}{u}.
%\end{equation*}
%Como $u - \sum_{i \in I} \inner{u}{e_i} e_i = 0$ (porque $(e_i)$ es una base), se sigue que:
%\begin{equation*}
%\norma{u}^2 = \sum_{i \in I} |\inner{u}{e_i}|^2.
%\end{equation*}
%\end{solution}

\section{Espacios de Hilbert separables}

\noindent
Ahora, unos ejercicios sobre espacios de Hilbert separables.
 $V$
Recordamos que un espacio $V$
 es separable si existe un subconjunto numerable $S \subset V$ denso en $V$. 

% --- Configuración de entornos ---
%\newtheorem{ex}{Ejercicio}
%\newenvironment{solution}{\begin{proof}[Solución]}{\end{proof}}

\begin{ex}
Demostrar que todo espacio de Hilbert de dimensión finita es separable.
\end{ex}

%\begin{solution}
%Sea $V$ un espacio de Hilbert de dimensión finita $n$. Entonces, existe una base ortonormal $(e_1, e_2, \ldots, e_n)$ de $V$. 

%El conjunto de todas las combinaciones lineales con coeficientes racionales de los vectores de la base:
%\begin{equation*}
%S = \left{ \sum_{i=1}^n q_i e_i : q_i \in \mathbb{Q} \right},
%\end{equation*}
%es numerable (porque $\mathbb{Q}^n$ es numerable) y denso en $V$ (porque $\mathbb{Q}$ es denso en $\mathbb{R}$). Por lo tanto, $V$ es separable.
%\end{solution}

\begin{ex}
¿Es el espacio $\ell^2(\mathbb{N})$ separable? Justificar la respuesta.
\end{ex}

%\begin{solution}
%Sí, $\ell^2(\mathbb{N})$ es separable. 
%
%Consideremos el conjunto $S$ de todas las sucesiones $(x_n)_{n \in \mathbb{N}}$ tales que $x_n \in \mathbb{Q}$ para todo $n$ y $x_n = 0$ para todo $n$ suficientemente grande. Este conjunto es numerable (porque es una unión numerable de conjuntos finitos). 

%Además, $S$ es denso en $\ell^2(\mathbb{N})$, ya que para cualquier sucesión $(y_n) \in \ell^2(\mathbb{N})$ y cualquier $\epsilon > 0$, podemos aproximar cada $y_n$ por un número racional $q_n$ de tal manera que $\sum_{n=1}^\infty |y_n - q_n|^2 < \epsilon^2$. 

%Por lo tanto, $\ell^2(\mathbb{N})$ es separable.
%\end{solution}



%\begin{ex}
%Sea $V$ un espacio de Hilbert separable. Demostrar que toda base ortonormal de $V$ es numerable.
%\end{ex}

%\begin{solution}
%Sea $(e_i)_{i \in I}$ una base ortonormal de $V$. Como $V$ es separable, existe un subconjunto numerable $S \subset V$ denso en $V$. 

%Para cada $s \in S$, definimos el conjunto:
%\begin{equation*}
%I_s = { i \in I : \inner{s}{e_i} \neq 0 }.
%\end{equation*}
%Como $s = \sum_{i \in I} \inner{s}{e_i} e_i$, el conjunto $I_s$ es numerable (porque la suma tiene a lo sumo una cantidad numerable de términos no nulos). 

%Ahora, observamos que:
%\begin{equation*}
%I = \bigcup_{s \in S} I_s,
%\end{equation*}
%ya que si $i \in I$, entonces $e_i \neq 0$ y existe $s \in S$ tal que $\inner{s}{e_i} \neq 0$ (porque $S$ es denso en $V$). 

%Como $S$ es numerable y cada $I_s$ es numerable, $I$ es numerable.
%\end{solution}

\begin{ex}
Demostrar que en un espacio de Hilbert separable, toda sucesión acotada tiene una subsucesión débilmente convergente.
\end{ex}

%\begin{solution}
%Sea $(u_n)$ una sucesión acotada en un espacio de Hilbert separable $V$. Como $V$ es separable, existe una base ortonormal numerable $(e_k)_{k \in \mathbb{N}}$. 

%Para cada $k \in \mathbb{N}$, la sucesión $(\inner{u_n}{e_k})$ está acotada en $\mathbb{R}$ (por la desigualdad de Cauchy-Schwarz). Por el teorema de Bolzano-Weierstrass, existe una subsucesión $(u_{n,1})$ tal que $(\inner{u_{n,1}}{e_1})$ converge. 

%Aplicando el mismo argumento a $(u_{n,1})$, obtenemos una subsucesión $(u_{n,2})$ tal que $(\inner{u_{n,2}}{e_2})$ converge. Continuando este proceso inductivamente, obtenemos una subsucesión diagonal $(u_{n,n})$ tal que $(\inner{u_{n,n}}{e_k})$ converge para todo $k \in \mathbb{N}$. 

%Por la Proposición
%ef{prop:hilbert-separable} (equivalencia entre convergencia débil y convergencia de componentes en una base hilbertiana), $(u_{n,n})$ converge débilmente en $V$.
%\end{solution}

%\begin{ex}
%Sea $V$ un espacio de Hilbert separable. ¿Es el espacio $L^2([0,1])$ separable? Justificar la respuesta.
%\end{ex}

%\begin{solution}
%Sí, $L^2([0,1])$ es separable. 
%
%El conjunto de las funciones polinómicas con coeficientes racionales es numerable y denso en $L^2([0,1])$ (por el teorema de aproximación de Weierstrass y la densidad de las funciones continuas en $L^2([0,1])$). Por lo tanto, $L^2([0,1])$ es separable.
%\end{solution}

%\begin{ex}
%Demostrar que si $V$ es un espacio de Hilbert separable, entonces toda base ortonormal de $V$ es completa.
%\end{ex}

%\begin{solution}
%Sea $(e_n)_{n \in \mathbb{N}}$ una base ortonormal de $V$. Supongamos que existe $u \in V$ tal que $\inner{u}{e_n} = 0$ para todo $n \in \mathbb{N}$. 

%Como $V$ es separable, existe un subconjunto numerable denso $S \subset V$. Para todo $s \in S$, podemos escribir $s$ como:
%\begin{equation*}
%s = \sum_{n=1}^\infty \inner{s}{e_n} e_n.
%\end{equation*}
%Si $\inner{u}{s} = 0$ para todo $s \in S$, entonces, por densidad de $S$, $\inner{u}{v} = 0$ para todo $v \in V$. En particular, $\inner{u}{u} = 0$, lo que implica $u = 0$. 

%Por lo tanto, $(e_n)$ es completa.
%\end{solution}



\begin{ex}
Sea $V$ un espacio de Hilbert separable y sea $(e_n){n \in \mathbb{N}}$ una base ortonormal de $V$. Demostrar que para todo $u \in V$ se tiene:
\begin{equation*}
\norma{u}^2 = \sum_{n=1}^\infty \inner{u}{e_n}^2.
\end{equation*}
\end{ex}

%\begin{solution}
%Por el teorema de Pitágoras en espacios de Hilbert, para cualquier $N \in \mathbb{N}$, se tiene:
%\begin{equation*}
%\norma{\sum_{n=1}^N \inner{u}{e_n} e_n}^2 = \sum_{n=1}^N |\inner{u}{e_n}|^2.
%\end{equation*}
%Tomando límite cuando $N \to \infty$, obtenemos:
%\begin{equation*}
%\norma{u}^2 = \norma{\sum_{n=1}^\infty \inner{u}{e_n} e_n}^2 = \sum_{n=1}^\infty |\inner{u}{e_n}|^2,
%\end{equation*}
%ya que $u = \sum_{n=1}^\infty \inner{u}{e_n} e_n$ (porque $(e_n)$ es una base ortonormal).
%\end{solution}

%\begin{ex}
%Demostrar que en un espacio de Hilbert separable, toda sucesión débilmente convergente está acotada.
%\end{ex}

%\begin{solution}
%Sea $(u_n)$ una sucesión en $V$ tal que $u_n \rightharpoonup u$. Por la definición de convergencia débil, para todo $v \in V$, la sucesión $(\inner{u_n}{v})$ está acotada. 

%En particular, para $v = u_n$, se tiene:
%\begin{equation*}
%|\inner{u_n}{u_n}| = \norma{u_n}^2 \leq M_v,
%\end{equation*}
%donde $M_v$ es una constante que depende de $v$. Sin embargo, esto no es suficiente para concluir directamente. 

%En su lugar, usamos el principio de acotación uniforme: si $(u_n)$ no estuviera acotada, existiría una subsucesión $(u_{n_k})$ tal que $\norma{u_{n_k}} \to \infty$. Pero entonces, para $v = u$, tendríamos:
%\begin{equation*}
%|\inner{u_{n_k}}{u}| \leq \norma{u_{n_k}} \norma{u} \to \infty,
%\end{equation*}
%lo que contradice que $(\inner{u_{n_k}}{u})$ converge. Por lo tanto, $(u_n)$ está acotada.
%\end{solution}

%\begin{ex}
%Sea $V$ un espacio de Hilbert separable y sea $(e_n)_{n \in \mathbb{N}}$ una base ortonormal. Demostrar que el conjunto ${e_n : n \in \mathbb{N}}$ es cerrado en $V$.
%\end{ex}

%\begin{solution}
%El conjunto ${e_n : n \in \mathbb{N}}$ no es cerrado en $V$ en general. Por ejemplo, en $\ell^2(\mathbb{N})$, la sucesión $(e_n)$ converge débilmente a $0$, pero $0 \notin {e_n : n \in \mathbb{N}}$. 

%Sin embargo, el conjunto ${e_n : n \in \mathbb{N}}$ es débilmente cerrado en el sentido de que si una sucesión en ${e_n : n \in \mathbb{N}}$ converge débilmente a un punto en $V$, entonces ese punto debe ser uno de los $e_n$. 

%Pero en la topología fuerte, el conjunto no es cerrado. 

%Corrección del enunciado: Si el ejercicio pretendía preguntar por la clausura débil, entonces la respuesta es que el conjunto ${e_n : n \in \mathbb{N}}$ es débilmente cerrado. 

%Para ver esto, supongamos que $(e_{n_k})$ converge débilmente a $u \in V$. Entonces, para todo $m \in \mathbb{N}$:
%\begin{equation*}
%\inner{e_{n_k}}{e_m} \to \inner{u}{e_m}.
%\end{equation*}
%Como $\inner{e_{n_k}}{e_m} = \delta_{n_k m}$ (donde $\delta_{n_k m}$ es la delta de Kronecker), se tiene que $\inner{u}{e_m} = 1$ si $m$ está en la subsucesión $(n_k)$ y $0$ en caso contrario. Por lo tanto, $u = e_m$ para algún $m$.
%\end{solution}

%\begin{ex}
%Demostrar que si $V$ es un espacio de Hilbert separable, entonces $V$ tiene una base ortonormal numerable.
%\end{ex}

%\begin{solution}
%Como $V$ es separable, existe un subconjunto numerable $S = {s_1, s_2, \ldots} \subset V$ denso en $V$. 

%Aplicamos el proceso de ortonormalización de Gram-Schmidt a $S$ para obtener una familia ortonormal $(e_n)_{n \in \mathbb{N}}$. 


%Ortonormalidad: Por construcción, $(e_n)$ es ortonormal.

%Base: Para todo $v \in V$ y $\epsilon > 0$, existe $s \in S$ tal que $\norma{v - s} < \epsilon$. Como $s$ es combinación lineal finita de los $e_n$, se tiene que $(e_n)$ es total en $V$.

%Por lo tanto, $(e_n)$ es una base ortonormal numerable de $V$.
%\end{solution}

\section{Espacio de Lebesgue}

\noindent
Ahora, algunos ejercicios sobre el espacio $L^2$. % en dimensión 1

\begin{ex}
Demostrar que el espacio $\Ldos([a, b])$ es un espacio vectorial sobre $\mathbb{R}$.
\end{ex}

%\begin{solution}
%Para demostrar que $\Ldos([a, b])$ es un espacio vectorial, debemos verificar las siguientes propiedades para funciones $f, g \in \Ldos([a, b])$ y $\alpha, \beta \in \mathbb{R}$:

%Cierre bajo suma:
% Por la desigualdad de Minkowski para integrales, se tiene:
% \begin{equation*}
% \norm{f + g}^2 = \int_a^b |f(x) + g(x)|^2 , dx \leq \left( \sqrt{\int_a^b |f(x)|^2 , dx} + \sqrt{\int_a^b |g(x)|^2 , dx} \right)^2 < \infty.
% \end{equation*}
% Por lo tanto, $f + g \in \Ldos([a, b])$.

%Cierre bajo multiplicación por escalar:
% \begin{equation*}
% \norm{\alpha f}^2 = \int_a^b |\alpha f(x)|^2 , dx = |\alpha|^2 \int_a^b |f(x)|^2 , dx < \infty.
% \end{equation*}
% Por lo tanto, $\alpha f \in \Ldos([a, b])$.

%Existencia del elemento neutro: La función nula $0(x) = 0$ pertenece a $\Ldos([a, b])$ y satisface $f + 0 = f$ para todo $f \in \Ldos([a, b])$.

%Existencia del opuesto: Para cada $f \in \Ldos([a, b])$, la función $-f$ también pertenece a $\Ldos([a, b])$ y satisface $f + (-f) = 0$.

%Por lo tanto, $\Ldos([a, b])$ es un espacio vectorial.
%\end{solution}

\begin{ex}
Demostrar que la función $f(x) = x$ pertenece a $\Ldos([0, 1])$.
\end{ex}

%\begin{solution}
%Calculamos la integral de $|f(x)|^2$ en $[0, 1]$:
%\begin{equation*}
%\int_0^1 |x|^2 , dx = \int_0^1 x^2 , dx = \left[ \frac{x^3}{3} \right]_0^1 = \frac{1}{3} < \infty.
%\end{equation*}
%Por lo tanto, $f \in \Ldos([0, 1])$.
%\end{solution}

\begin{ex}
Demostrar que la función $f(x) = {1}/{x}$ no pertenece a $\Ldos([0, 1])$.
\end{ex}

%\begin{solution}
%Calculamos la integral de $|f(x)|^2$ en $[0, 1]$:
%\begin{equation*}
%\int_0^1 \left| \frac{1}{x} \right|^2 , dx = \int_0^1 \frac{1}{x^2} , dx = \left[ -\frac{1}{x} \right]_0^1 = +\infty.
%\end{equation*}
%Por lo tanto, $f \notin \Ldos([0, 1])$.
%\end{solution}

\begin{ex}
Sea $f \in \Ldos([a, b])$. Demostrar que para cualquier $\alpha \in \mathbb{R}$, la función $g(x) = f(x + \alpha)$ (traslación) también pertenece a $\Ldos([a - \alpha, b - \alpha])$.
\end{ex}

%\begin{solution}
%Hacemos un cambio de variable $y = x + \alpha$. Entonces, $dy = dx$, y los límites de integración cambian a $y \in [a, b]$ cuando $x \in [a - \alpha, b - \alpha]$. Por lo tanto:
%\begin{align*}
%\int_{a - \alpha}^{b - \alpha} |g(x)|^2 , dx &= \int_{a - \alpha}^{b - \alpha} |f(x + \alpha)|^2 , dx \
%&= \int_a^b |f(y)|^2 , dy \
%&= \norm{f}^2 < \infty.
%\end{align*}
%Por lo tanto, $g \in \Ldos([a - \alpha, b - \alpha])$.
%\end{solution}

%\begin{ex}
%Demostrar que el espacio $\Ldos([a, b])$ es completo, es decir, es un espacio de Hilbert.
%\end{ex}

%\begin{solution}
%Para demostrar que $\Ldos([a, b])$ es completo, debemos mostrar que toda sucesión de Cauchy en $\Ldos([a, b])$ converge a una función en $\Ldos([a, b])$. 

%Sea $(f_n)$ una sucesión de Cauchy en $\Ldos([a, b])$. Entonces, para todo $\epsilon > 0$, existe $N \in \mathbb{N}$ tal que para todo $m, n \geq N$, se tiene:
%\begin{equation*}
%\norm{f_n - f_m} < \epsilon.
%\end{equation*}

%Por el teorema de Riesz-Fischer, existe una subsucesión $(f_{n_k})$ y una función $f \in \Ldos([a, b])$ tal que $f_{n_k} \to f$ casi todo punto y en la norma de $\Ldos([a, b])$. 

%Como $(f_n)$ es de Cauchy, se puede mostrar que toda la sucesión $(f_n)$ converge a $f$ en la norma de $\Ldos([a, b])$. Por lo tanto, $\Ldos([a, b])$ es completo.
%\end{solution}

%\begin{ex}
%Demostrar que el conjunto de las funciones continuas en $[a, b]$, denotado por $C([a, b])$, es denso en $\Ldos([a, b])$.
%\end{ex}

%\begin{solution}
%Por el teorema de aproximación de Weierstrass, toda función continua en $[a, b]$ puede aproximarse uniformemente por polinomios. Además, toda función en $\Ldos([a, b])$ puede aproximarse por funciones continuas en la norma de $\Ldos([a, b])$. 

%Más precisamente, para cualquier $f \in \Ldos([a, b])$ y $\epsilon > 0$, existe una función continua $g \in C([a, b])$ tal que $\norm{f - g} < \epsilon$. Esto se sigue del hecho de que las funciones escalonadas son densas en $\Ldos([a, b])$ y las funciones continuas pueden aproximar a las funciones escalonadas arbitrariamente bien.
%\end{solution}

\begin{ex}
Sea $f \in \Ldos([0, 1])$. Demostrar que la función $g(x) = f(x) \cdot x$ también pertenece a $\Ldos([0, 1])$.
\end{ex}

%\begin{solution}
%Calculamos la norma de $g$:
%\begin{align*}
%\norm{g}^2 &= \int_0^1 |f(x) x|^2 , dx \
%&= \int_0^1 |f(x)|^2 x^2 , dx \
%&\leq \int_0^1 |f(x)|^2 , dx \quad \text{(ya que } x^2 \leq 1 \text{ en } [0, 1]) \
%&= \norm{f}^2 < \infty.
%\end{align*}
%Por lo tanto, $g \in \Ldos([0, 1])$.
%\end{solution}

%\begin{ex}
%Demostrar que el producto interno en $\Ldos([a, b])$ está bien definido y satisface:
%\begin{equation*}
%\inner{f}{g} = \int_a^b f(x) g(x) , dx \quad \forall f, g \in \Ldos([a, b]).
%\end{equation*}
%\end{ex}

%\begin{solution}
%Para demostrar que el producto interno está bien definido, debemos mostrar que la integral $\int_a^b f(x) g(x) , dx$ existe y es finita para todo $f, g \in \Ldos([a, b])$. 

%Por la desigualdad de Cauchy-Schwarz para integrales:
%\begin{equation*}
%\left| \int_a^b f(x) g(x) , dx \right| \leq \int_a^b |f(x) g(x)| , dx \leq \sqrt{\int_a^b |f(x)|^2 , dx} \sqrt{\int_a^b |g(x)|^2 , dx} = \norm{f} \norm{g} < \infty.
%\end{equation*}
%Por lo tanto, el producto interno está bien definido. Además, es claro que satisface las propiedades de un producto interno (linealidad, simetría y definitud positiva).
%\end{solution}

%\begin{ex}
%Demostrar que si $f \in \Ldos([a, b])$ y $f(x) \geq 0$ casi todo punto, entonces:
%\begin{equation*}
%\norm{f}^2 = \int_a^b f(x)^2 , dx \geq 0.
%\end{equation*}
%\end{ex}

%\begin{solution}
%Como $f(x) \geq 0$ casi todo punto, se tiene que $f(x)^2 \geq 0$ casi todo punto. Por lo tanto, la integral de una función no negativa es no negativa:
%\begin{equation*}
%\int_a^b f(x)^2 , dx \geq 0.
%\end{equation*}
%Además, $\norm{f}^2 = 0$ si y solo si $f(x) = 0$ casi todo punto, lo que es consistente con la definición de la norma en $\Ldos([a, b])$.
%\end{solution}


%\begin{ex}
%Demostrar que el conjunto de las funciones polinómicas en $[a, b]$ es denso en $\Ldos([a, b])$.
%\end{ex}

%\begin{solution}
%Por el teorema de Weierstrass, toda función continua en $[a, b]$ puede aproximarse uniformemente por polinomios. Además, como las funciones continuas son densas en $\Ldos([a, b])$, se sigue que los polinomios también son densos en $\Ldos([a, b])$. 

%Más formalmente, para cualquier $f \in \Ldos([a, b])$ y $\epsilon > 0$, existe una función continua $g \in C([a, b])$ tal que $\norm{f - g} < \epsilon/2$. Luego, existe un polinomio $p$ tal que $\norm{g - p}_\infty < \epsilon/(2 \sqrt{b - a})$. Por lo tanto:
%\begin{align*}
%\norm{f - p} &= \norm{f - g + g - p} \
%&\leq \norm{f - g} + \norm{g - p} \
%&< \frac{\epsilon}{2} + \frac{\epsilon}{2 \sqrt{b - a}} \sqrt{b - a} \
%&= \epsilon.
%\end{align*}
%Por lo tanto, los polinomios son densos en $\Ldos([a, b])$.
%\end{solution}

%\begin{ex}
%Demostrar que si $f \in \Ldos([a, b])$ y $g \in \L^\infty([a, b])$, entonces $f \cdot g \in \Ldos([a, b])$.
%\end{ex}

%\begin{solution}
%Como $g \in \L^\infty([a, b])$, existe $M > 0$ tal que $|g(x)| \leq M$ casi todo punto. Por lo tanto:
%\begin{align*}
%\norm{f \cdot g}^2 &= \int_a^b |f(x) g(x)|^2 , dx \
%&= \int_a^b |f(x)|^2 |g(x)|^2 , dx \
%&\leq M^2 \int_a^b |f(x)|^2 , dx \
%&= M^2 \norm{f}^2 < \infty.
%\end{align*}
%Por lo tanto, $f \cdot g \in \Ldos([a, b])$.
%\end{solution}

%%%%%%%%%%%%%%%%%%%%%%%%%%%%%%%%%%%%%%%%%%%%%%%%%%

\section{Espacios de Sobolev}

\noindent % ---
En lo siguiente, varios ejercicios sobre espacios de Sobolev.
% ---

\begin{ex}
Dar la definición del espacio de Sobolev $\Huno((a, b))$ en dimensión 1.
\end{ex}

%\begin{solution}
%El espacio de Sobolev $\Huno((a, b))$ se define como:
%\begin{equation*}
%\Huno((a, b)) = \left{ u \in \Ldos((a, b)) : u' \in \Ldos((a, b)) \right},
%\end{equation*}
%donde $u'$ es la derivada débil de $u$. 

%La norma en $\Huno((a, b))$ está dada por:
%\begin{equation*}
%\norm{u}_{\Huno((a, b))} = \left( \int_a^b |u(x)|^2 , dx + \int_a^b |u'(x)|^2 , dx \right)^{1/2}.
%\end{equation*}
%\end{solution}

%\begin{ex}
%Demostrar que $\Huno((a, b))$ es un espacio de Hilbert.
%\end{ex}

%\begin{solution}
%Para demostrar que $\Huno((a, b))$ es un espacio de Hilbert, debemos mostrar que es un espacio vectorial con un producto interno que lo hace completo.


%Espacio vectorial:


%Suma: Si $u, v \in \Huno((a, b))$, entonces $u + v \in \Ldos((a, b))$ y $(u + v)' = u' + v' \in \Ldos((a, b))$, por lo que $u + v \in \Huno((a, b))$.

%Multiplicación por escalar: Si $u \in \Huno((a, b))$ y $\alpha \in \mathbb{R}$, entonces $\alpha u \in \Ldos((a, b))$ y $(\alpha u)' = \alpha u' \in \Ldos((a, b))$, por lo que $\alpha u \in \Huno((a, b))$.


%Producto interno: Definimos el producto interno en $\Huno((a, b))$ como:
% \begin{equation*}
% \inner{u}{v}{\Huno} = \int_a^b u(x) v(x) , dx + \int_a^b u'(x) v'(x) , dx.
% \end{equation*}
% Este producto interno induce la norma $\norm{u}{\Huno} = \sqrt{\inner{u}{u}_{\Huno}}$.


%Completitud:
% Sea $(u_n)$ una sucesión de Cauchy en $\Huno((a, b))$. Entonces, $(u_n)$ y $(u_n')$ son sucesiones de Cauchy en $\Ldos((a, b))$. Como $\Ldos((a, b))$ es completo, existen $u, v \in \Ldos((a, b))$ tales que $u_n \to u$ y $u_n' \to v$ en $\Ldos((a, b))$.

%Además, como la derivada débil es cerrada, se tiene que $v = u'$. Por lo tanto, $u \in \Huno((a, b))$ y $u_n \to u$ en $\Huno((a, b))$. 

% En consecuencia, $\Huno((a, b))$ es completo y, por lo tanto, un espacio de Hilbert.

%\end{solution}



\begin{ex}
Demostrar que toda función en $\Huno((a, b))$ es continua en $[a, b]$.
\end{ex}

%\begin{solution}
%Sea $u \in \Huno((a, b))$. Por definición, $u \in \Ldos((a, b))$ y $u' \in \Ldos((a, b))$. 

%Para cualquier $x, y \in [a, b]$, se tiene:
%\begin{equation*}
%u(x) - u(y) = \int_y^x u'(t) , dt.
%\end{equation*}
%Por la desigualdad de Cauchy-Schwarz:
%\begin{align*}
%|u(x) - u(y)| &= \left| \int_y^x u'(t) , dt \right| \
%&\leq \int_y^x |u'(t)| , dt \
%&\leq |x - y|^{1/2} \left( \int_y^x |u'(t)|^2 , dt \right)^{1/2} \
%&\leq |x - y|^{1/2} \norm{u'}_{\Ldos((a, b))}.
%\end{align*}
%Esto muestra que $u$ es Hölder-continua con exponente $1/2$. En particular, $u$ es continua en $[a, b]$.
%\end{solution}

\begin{ex}
Demostrar que el espacio $C^1([a, b])$ (funciones continuamente diferenciables en $[a, b]$) está contenido en $\Huno((a, b))$.
\end{ex}

%\begin{solution}
%Sea $u \in C^1([a, b])$. Entonces, $u$ es continua en $[a, b]$ y su derivada $u'$ también es continua en $[a, b]$. 

%Como $[a, b]$ es un intervalo acotado, $u$ y $u'$ están acotadas en $[a, b]$. Por lo tanto:
%\begin{align*}
%\int_a^b |u(x)|^2 , dx &\leq (b - a) \max_{x \in [a, b]} |u(x)|^2 < \infty, \
%\int_a^b |u'(x)|^2 , dx &\leq (b - a) \max_{x \in [a, b]} |u'(x)|^2 < \infty.
%\end{align*}
%Por lo tanto, $u \in \Ldos((a, b))$ y $u' \in \Ldos((a, b))$, lo que implica que $u \in \Huno((a, b))$.
%\end{solution}

%\begin{ex}
%Demostrar que el espacio $C^\infty_0((a, b))$ (funciones infinitamente diferenciables con soporte compacto en $(a, b)$) es denso en $\Huno_0((a, b))$.
%\end{ex}

%\begin{solution}
%El espacio $\Huno_0((a, b))$ se define como el cierre de $C^\infty_0((a, b))$ en $\Huno((a, b))$. 

%Para demostrar que $C^\infty_0((a, b))$ es denso en $\Huno_0((a, b))$, consideramos una función $u \in \Huno_0((a, b))$. Por definición, existe una sucesión $(u_n) \subset C^\infty_0((a, b))$ tal que $u_n \to u$ en $\Huno((a, b))$. 

%Esto significa que:
%\begin{equation*}
%\norm{u_n - u}_{\Huno} = \left( \int_a^b |u_n(x) - u(x)|^2 , dx + \int_a^b |u_n'(x) - u'(x)|^2 , dx \right)^{1/2} \to 0.
%\end{equation*}
%Por lo tanto, $C^\infty_0((a, b))$ es denso en $\Huno_0((a, b))$.
%\end{solution}

\begin{ex}
Demostrar que en $\Huno_0((a, b))$, la norma de $u$ en ${\Huno}$ es equivalente a la norma de $u'$ en ${\Ldos}$.
\end{ex}

%\begin{solution}
%Para demostrar la equivalencia de las normas, debemos mostrar que existen constantes $C_1, C_2 > 0$ tales que:
%\begin{equation*}
%C_1 \norm{u'}{\Ldos} \leq \norm{u}{\Huno} \leq C_2 \norm{u'}_{\Ldos} \quad \forall u \in \Huno_0((a, b)).
%\end{equation*}

%Primera desigualdad:
% Es claro que $\norm{u'}{\Ldos} \leq \norm{u}{\Huno}$, por lo que podemos tomar $C_1 = 1$.

%Segunda desigualdad (Desigualdad de Poincaré en dimensión 1):
% Sea $u \in \Huno_0((a, b))$. Como $u(a) = u(b) = 0$ (porque $u \in \Huno_0((a, b))$), para cualquier $x \in [a, b]$, se tiene:
% \begin{equation*}
%u(x) = \int_a^x u'(t) , dt.
% \end{equation*}
% Por la desigualdad de Cauchy-Schwarz:
% \begin{align*}
% |u(x)|^2 &= \left| \int_a^x u'(t) , dt \right|^2 \
% &\leq (x - a) \int_a^x |u'(t)|^2 , dt \
% &\leq (b - a) \int_a^b |u'(t)|^2 , dt \
% &= (b - a) \norm{u'}^2_{\Ldos}.
% \end{align*}
% Integrando sobre $[a, b]$:
% \begin{align*}
% \int_a^b |u(x)|^2 , dx &\leq (b - a)^2 \int_a^b |u'(x)|^2 , dx \
% &= (b - a)^2 \norm{u'}^2_{\Ldos}.
% \end{align*}
% Por lo tanto:
% \begin{equation*}
% \norm{u}^2_{\Huno} = \norm{u}^2_{\Ldos} + \norm{u'}^2_{\Ldos} \leq (1 + (b - a)^2) \norm{u'}^2_{\Ldos}.
% \end{equation*}
% Tomando $C_2 = \sqrt{1 + (b - a)^2}$, obtenemos la desigualdad deseada.
%\end{solution}



\begin{ex}
Demostrar que si $u \in \Huno((a, b))$, entonces $u$ es absolutamente continua en $[a, b]$. \emph{[Una función $u$ es absolutamente continua en $[a, b]$ si para todo $\epsilon > 0$, existe $\delta > 0$ tal que para cualquier colección finita de intervalos disjuntos $(x_i, y_i) \subset [a, b]$ con $\sum_i |y_i - x_i| < \delta$, se tiene $\sum_i |u(y_i) - u(x_i)| < \epsilon$.]}
\end{ex}

%\begin{solution}
%
%Como $u \in \Huno((a, b))$, su derivada débil $u' \in \Ldos((a, b))$. Por el teorema fundamental del cálculo para funciones absolutamente continuas, si $u'$ existe casi todo punto y es integrable, entonces $u$ es absolutamente continua. 
%
%En particular, para cualquier $x, y \in [a, b]$, se tiene:
%\begin{equation*}
%u(y) - u(x) = \int_x^y u'(t) , dt,
%\end{equation*}
%y por la desigualdad de Hölder:
%\begin{equation*}
%|u(y) - u(x)| \leq \int_x^y |u'(t)| , dt \leq |y - x|^{1/2} \left( \int_x^y |u'(t)|^2 , dt \right)^{1/2} \leq |y - x|^{1/2} \norm{u'}_{\Ldos}.
%\end{equation*}
%Esto muestra que $u$ es Hölder-continua con exponente $1/2$, y por lo tanto, absolutamente continua.
%\end{solution}


\begin{ex}
Demostrar que el espacio $\Huno((a, b))$ está contenido en $C([a, b])$ y que la inclusión es continua.
\end{ex}

%\begin{solution}

%Inclusión en $C([a, b])$:
% Como se demostró en el Ejercicio 3, toda función en $\Huno((a, b))$ es continua en $[a, b]$. Por lo tanto, $\Huno((a, b)) \subset C([a, b])$.

%Continuidad de la inclusión:
% Debemos mostrar que existe una constante $C > 0$ tal que:
% \begin{equation*}
% \norm{u}{C([a, b])} \leq C \norm{u}{\Huno} \quad \forall u \in \Huno((a, b)).
% \end{equation*}
% Para cualquier $x \in [a, b]$, se tiene:
% \begin{align*}
% |u(x)| &= \left| u(a) + \int_a^x u'(t) , dt \right| \
% &\leq |u(a)| + \int_a^x |u'(t)| , dt \
% &\leq |u(a)| + |x - a|^{1/2} \left( \int_a^x |u'(t)|^2 , dt \right)^{1/2} \
% &\leq |u(a)| + (b - a)^{1/2} \norm{u'}{\Ldos}.
% \end{align*}
% Además, por la desigualdad de Cauchy-Schwarz:
% \begin{equation*}
% |u(a)| \leq \frac{1}{b - a} \int_a^b |u(x)| , dx \leq \frac{1}{\sqrt{b - a}} \norm{u}{\Ldos}.
% \end{equation*}
% Por lo tanto:
% \begin{align*}
% \norm{u}{C([a, b])} &= \max{x \in [a, b]} |u(x)| \
% &\leq \frac{1}{\sqrt{b - a}} \norm{u}{\Ldos} + (b - a)^{1/2} \norm{u'}{\Ldos} \
% &\leq C \left( \norm{u}{\Ldos} + \norm{u'}{\Ldos} \right) \
% &\leq C \norm{u}_{\Huno},
% \end{align*}
% donde $C = \max \left{ \frac{1}{\sqrt{b - a}}, (b - a)^{1/2} \right}$.

%Por lo tanto, la inclusión es continua.

%\end{solution}

%\begin{ex}
%Demostrar que si $u \in \Huno((a, b))$ y $u' = 0$ casi todo punto, entonces $u$ es constante en $[a, b]$.
%\end{ex}

%\begin{solution}
%Como $u \in \Huno((a, b))$, $u$ es continua en $[a, b]$ (Ejercicio 3). Además, $u' = 0$ casi todo punto. 

%Para cualquier $x, y \in [a, b]$, se tiene:
%\begin{equation*}
%u(y) - u(x) = \int_x^y u'(t) , dt = 0,
%\end{equation*}
%ya que $u' = 0$ casi todo punto. Por lo tanto, $u(y) = u(x)$ para todo $x, y \in [a, b]$, lo que implica que $u$ es constante en $[a, b]$.
%\end{solution}

%\begin{ex}
%Demostrar que el espacio $\Huno((a, b))$ es separable.
%\end{ex}

%\begin{solution}
%Para demostrar que $\Huno((a, b))$ es separable, utilizamos el hecho de que $\Ldos((a, b))$ es separable y que $C^\infty_0((a, b))$ es denso en $\Huno((a, b))$.

%Separabilidad de $\Ldos((a, b))$:
% Como se vio en los ejercicios anteriores, $\Ldos((a, b))$ es separable.

%Densidad de $C^\infty_0((a, b))$ en $\Huno((a, b))$:
% El espacio $C^\infty_0((a, b))$ es numerable y denso en $\Huno((a, b))$ (Ejercicio 4).

%Separabilidad de $\Huno((a, b))$:
% Como $\Huno((a, b))$ contiene un subconjunto numerable denso ($C^\infty_0((a, b))$), es separable.
%\end{solution}

%\begin{ex}
%Demostrar que si $u \in \Huno((a, b))$ y $u \geq 0$ casi todo punto, entonces $u$ es no negativa en todo $[a, b]$.
%\end{ex}

%\begin{solution}
%Como $u \in \Huno((a, b))$, $u$ es continua en $[a, b]$ (Ejercicio 3). Además, $u \geq 0$ casi todo punto. 

%Supongamos por contradicción que existe $x_0 \in [a, b]$ tal que $u(x_0) < 0$. Como $u$ es continua, existe un intervalo $(x_0 - \delta, x_0 + \delta) \subset [a, b]$ tal que $u(x) < 0$ para todo $x \in (x_0 - \delta, x_0 + \delta)$. 

%Sin embargo, esto contradice que $u \geq 0$ casi todo punto, ya que $(x_0 - \delta, x_0 + \delta)$ tiene medida positiva. Por lo tanto, $u(x) \geq 0$ para todo $x \in [a, b]$.
%\end{solution}

%\begin{ex}
%Demostrar que el espacio $\Huno_0((a, b))$ es un subespacio cerrado de $\Huno((a, b))$.
%\end{ex}

%\begin{solution}
%El espacio $\Huno_0((a, b))$ se define como el cierre de $C^\infty_0((a, b))$ en $\Huno((a, b))$. Para demostrar que es un subespacio cerrado, debemos mostrar que si $(u_n) \subset \Huno_0((a, b))$ y $u_n \to u$ en $\Huno((a, b))$, entonces $u \in \Huno_0((a, b))$. 

%Como $u_n \in \Huno_0((a, b))$, se tiene que $u_n(a) = u_n(b) = 0$ para todo $n$. Además, por la continuidad de la inclusión de $\Huno((a, b))$ en $C([a, b])$ (Ejercicio 6), se tiene que $u_n \to u$ uniformemente en $[a, b]$. 

%Por lo tanto:
%\begin{equation*}
%u(a) = \lim_{n \to \infty} u_n(a) = 0, \quad u(b) = \lim_{n \to \infty} u_n(b) = 0.
%\end{equation*}
%Esto implica que $u \in \Huno_0((a, b))$, y por lo tanto, $\Huno_0((a, b))$ es un subespacio cerrado de $\Huno((a, b))$.
%\end{solution}

\section{Calculo diferencial}

% ---
% Ejercicios sobre cálculo diferencial en espacios de Hilbert (Diferencial de Fréchet)
% ---

\noindent
Finalmente una serie sobre cálculo diferencial en espacios de Hilbert (diferencial de Fréchet).

\begin{ex}
Dar la definición de diferencial de Fréchet para una función $F: U \subset V \to W$, donde $V$ y $W$ son espacios de Hilbert.
\end{ex}

%\begin{solution}
%Sea $F: U \subset V \to W$ una función definida en un abierto $U$ de un espacio de Hilbert $V$, con valores en otro espacio de Hilbert $W$. Decimos que $F$ es diferenciable en el sentido de Fréchet en un punto $u \in U$ si existe una aplicación lineal y continua $DF(u): V \to W$ tal que:
%\begin{equation*}
%\lim_{\norma{h}_V \to 0} \frac{\norma{F(u + h) - F(u) - DF(u) h}_W}{\norma{h}_V} = 0.
%\end{equation*}
%La aplicación lineal $DF(u)$ se conoce como la diferencial de Fréchet de $F$ en $u$.
%\end{solution}

\begin{ex}
Demostrar que si $F: V \to W$ es una aplicación lineal y continua entre espacios de Hilbert, entonces $F$ es diferenciable en el sentido de Fréchet en todo punto $u \in V$, y su diferencial es $DF(u) = F$ para todo $u \in V$.
\end{ex}

%\begin{solution}
%Para cualquier $u, h \in V$, se tiene:
%\begin{align*}
%F(u + h) - F(u) - F(h) &= F(u) + F(h) - F(u) - F(h) = 0.
%\end{align*}
%Por lo tanto:
%\begin{equation*}
%\lim_{\norma{h}_V \to 0} \frac{\norma{F(u + h) - F(u) - F(h)}_W}{\norma{h}V} = \lim{\norma{h}_V \to 0} \frac{\norma{0}_W}{\norma{h}_V} = 0.
%\end{equation*}
%Esto muestra que $F$ es diferenciable en el sentido de Fréchet en todo punto $u \in V$, y su diferencial es $DF(u) = F$.
%\end{solution}

\begin{ex}
Sea $V$ un espacio de Hilbert y consideremos la función $F: V \to \mathbb{R}$ definida por:
\begin{equation*}
F(u) = \norma{u}^2 = \inner{u}{u}.
\end{equation*}
Demostrar que $F$ es diferenciable en el sentido de Fréchet en todo punto $u \in V$ y encontrar su diferencial.
\end{ex}

%\begin{solution}
%Calculamos la diferencia:
%\begin{align*}
%F(u + h) - F(u) &= \norma{u + h}^2 - \norma{u}^2 \
%&= \inner{u + h}{u + h} - \inner{u}{u} \
%&= \inner{u}{u} + 2 \inner{u}{h} + \inner{h}{h} - \inner{u}{u} \
%&= 2 \inner{u}{h} + \norma{h}^2.
%\end{align*}
%Por lo tanto:
%\begin{align*}
%F(u + h) - F(u) - 2 \inner{u}{h} &= \norma{h}^2.
%\end{align*}
%Dividiendo por $\norma{h}$:
%\begin{equation*}
%\frac{\norma{F(u + h) - F(u) - 2 \inner{u}{h}}}{\norma{h}} = \frac{\norma{h}^2}{\norma{h}} = \norma{h} \to 0 \quad \text{cuando } \norma{h} \to 0.
%\end{equation*}
%Esto muestra que $F$ es diferenciable en el sentido de Fréchet en $u$, y su diferencial es:
%\begin{equation*}
%DF(u) h = 2 \inner{u}{h}.
%\end{equation*}
%\end{solution}

%\begin{ex}
%Sea $V$ un espacio de Hilbert y consideremos la función $F: V \to \mathbb{R}$ definida por:
%\begin{equation*}
%F(u) = \inner{u}{v},
%\end{equation*}
%para algún $v \in V$ fijo. Demostrar que $F$ es diferenciable en el sentido de Fréchet en todo punto $u \in V$ y encontrar su diferencial.
%\end{ex}

%\begin{solution}
%Como $F$ es un funcional lineal y continuo (por la desigualdad de Cauchy-Schwarz), por el Ejercicio 3, su diferencial de Fréchet en cualquier punto $u \in V$ es $DF(u) = F$. 

%Es decir:
%\begin{equation*}
%DF(u) h = \inner{h}{v} \quad \forall h \in V.
%\end{equation*}
%\end{solution}



%\begin{ex}
%Sea $V$ un espacio de Hilbert y consideremos la función $F: V \to V$ definida por:
%\begin{equation*}
%F(u) = u^2,
%\end{equation*}
%donde $u^2$ se interpreta como el producto de $u$ consigo mismo en el sentido del producto interno, es decir, $F(u) = \inner{u}{u} u$. 

%Demostrar que $F$ es diferenciable en el sentido de Fréchet en todo punto $u \in V$ y encontrar su diferencial.
%\end{ex}

%\begin{solution}
%Calculamos la diferencia:
%\begin{align*}
%F(u + h) - F(u) &= \inner{u + h}{u + h} (u + h) - \inner{u}{u} u \
%&= \left( \inner{u}{u} + 2 \inner{u}{h} + \inner{h}{h} \right) (u + h) - \inner{u}{u} u \
%&= \inner{u}{u} u + \inner{u}{u} h + 2 \inner{u}{h} u + 2 \inner{u}{h} h + \inner{h}{h} u + \inner{h}{h} h - \inner{u}{u} u \
%&= \inner{u}{u} h + 2 \inner{u}{h} u + \text{términos de orden superior}.
%\end{align*}
%Ignorando los términos de orden superior (aquellos que son $o(\norma{h})$), obtenemos:
%\begin{align*}
%F(u + h) - F(u) &= \inner{u}{u} h + 2 \inner{u}{h} u + o(\norma{h}).
%\end{align*}
%Por lo tanto, la diferencial de Fréchet de $F$ en $u$ es:
%\begin{equation*}
%DF(u) h = \inner{u}{u} h + 2 \inner{u}{h} u.
%\end{equation*}
%\end{solution}



%\begin{ex}
%Sea $F: V \to W$ una función diferenciable en el sentido de Fréchet en un punto $u \in V$, donde $V$ y $W$ son espacios de Hilbert. Demostrar que la diferencial de Fréchet $DF(u)$ es única.
%\end{ex}

%\begin{solution}
%Supongamos que existen dos aplicaciones lineales y continuas $L_1, L_2: V \to W$ tales que:
%\begin{align*}
%\lim_{\norma{h}_V \to 0} \frac{\norma{F(u + h) - F(u) - L_1 h}_W}{\norma{h}V} &= 0, \
%\lim{\norma{h}_V \to 0} \frac{\norma{F(u + h) - F(u) - L_2 h}_W}{\norma{h}_V} &= 0.
%\end{align*}
%Entonces, para cualquier $h \in V$:
%\begin{align*}
%\norma{L_1 h - L_2 h}_W &= \norma{(F(u + h) - F(u) - L_2 h) - (F(u + h) - F(u) - L_1 h)}_W \
%&\leq \norma{F(u + h) - F(u) - L_2 h}_W + \norma{F(u + h) - F(u) - L_1 h}_W.
%\end{align*}
%Dividiendo por $\norma{h}_V$ y tomando límite cuando $\norma{h}V \to 0$, obtenemos:
%\begin{equation*}
%\lim{\norma{h}_V \to 0} \frac{\norma{L_1 h - L_2 h}_W}{\norma{h}_V} = 0.
%\end{equation*}
%Como $L_1$ y $L_2$ son lineales y continuas, esto implica que $L_1 = L_2$.
%\end{solution}

\begin{ex}
Sea $F: V \to W$ una función constante, es decir, $F(u) = c$ para todo $u \in V$, donde $c \in W$ es fijo. Demostrar que $F$ es diferenciable en el sentido de Fréchet en todo punto $u \in V$ y encontrar su diferencial.
\end{ex}

%\begin{solution}
%Para cualquier $u, h \in V$, se tiene:
%\begin{equation*}
%F(u + h) - F(u) = c - c = 0.
%\end{equation*}
%Por lo tanto:
%\begin{equation*}
%\lim_{\norma{h}_V \to 0} \frac{\norma{F(u + h) - F(u) - 0}_W}{\norma{h}V} = \lim{\norma{h}_V \to 0} \frac{\norma{0}_W}{\norma{h}_V} = 0.
%\end{equation*}
%Esto muestra que $F$ es diferenciable en el sentido de Fréchet en todo punto $u \in V$, y su diferencial es la aplicación nula:
%\begin{equation*}
%DF(u) h = 0 \quad \forall h \in V.
%\end{equation*}
%\end{solution}

\begin{ex}
Sea $F: V \to W$ una función diferenciable en el sentido de Fréchet en un punto $u \in V$, y sea $G: W \to Z$ una función diferenciable en el sentido de Fréchet en $F(u)$, donde $V$, $W$ y $Z$ son espacios de Hilbert. Demostrar que la composición $G \circ F: V \to Z$ es diferenciable en el sentido de Fréchet en $u$ y encontrar su diferencial (regla de la cadena).
\end{ex}

%\begin{solution}
%Para demostrar la diferenciabilidad de $G \circ F$ en $u$, escribimos:
%\begin{align*}
%(G \circ F)(u + h) - (G \circ F)(u) &= G(F(u + h)) - G(F(u)) \
%&= G(F(u)) + DG(F(u)) (F(u + h) - F(u)) + o(\norma{F(u + h) - F(u)}_W) - G(F(u)) \
%&= DG(F(u)) (F(u + h) - F(u)) + o(\norma{F(u + h) - F(u)}_W) \
%&= DG(F(u)) (DF(u) h + o(\norma{h}_V)) + o(\norma{DF(u) h + o(\norma{h}_V)}_W) \
%&= DG(F(u)) DF(u) h + o(\norma{h}_V).
%\end{align*}
%Por lo tanto, la diferencial de Fréchet de $G \circ F$ en $u$ es:
%\begin{equation*}
%D(G \circ F)(u) = DG(F(u)) \circ DF(u).
%\end{equation*}
%\end{solution}

\begin{ex}
Sea $V$ un espacio de Hilbert y consideremos la función $F: V \to \mathbb{R}$ definida por:
\begin{equation*}
F(u) = \norma{u}.
\end{equation*}
Demostrar que $F$ es diferenciable en el sentido de Fréchet en todo punto $u \neq 0$ y encontrar su diferencial.
\end{ex}

%\begin{solution}
%Para $u \neq 0$, calculamos la diferencia:
%\begin{align*}
%F(u + h) - F(u) &= \norma{u + h} - \norma{u} \
%&= \sqrt{\inner{u + h}{u + h}} - \sqrt{\inner{u}{u}} \
%&= \sqrt{\norma{u}^2 + 2 \inner{u}{h} + \norma{h}^2} - \norma{u} \
%&= \norma{u} \sqrt{1 + \frac{2 \inner{u}{h}}{\norma{u}^2} + \frac{\norma{h}^2}{\norma{u}^2}} - \norma{u} \
%&= \norma{u} \left( 1 + \frac{\inner{u}{h}}{\norma{u}^2} + o(\norma{h}) - 1 \right) \
%&= \frac{\inner{u}{h}}{\norma{u}} + o(\norma{h}).
%\end{align*}
%Por lo tanto, la diferencial de Fréchet de $F$ en $u \neq 0$ es:
%\begin{equation*}
%DF(u) h = \frac{\inner{u}{h}}{\norma{u}}.
%\end{equation*}
%\end{solution}

\begin{ex}
Sea $F: V \to \mathbb{R}$ un funcional diferenciable en el sentido de Fréchet en un punto $u \in V$, donde $V$ es un espacio de Hilbert. Demostrar que si $F$ tiene un mínimo local en $u$, entonces $DF(u) = 0$.
\end{ex}

%\begin{solution}
%Supongamos que $F$ tiene un mínimo local en $u$. Entonces, existe $\delta > 0$ tal que para todo $h \in V$ con $\norma{h} < \delta$, se tiene:
%\begin{equation*}
%F(u + h) \geq F(u).
%\end{equation*}
%Como $F$ es diferenciable en $u$, podemos escribir:
%\begin{equation*}
%F(u + h) = F(u) + DF(u) h + o(\norma{h}).
%\end{equation*}
%Por lo tanto:
%\begin{equation*}
%F(u) + DF(u) h + o(\norma{h}) \geq F(u) \quad \Rightarrow \quad DF(u) h + o(\norma{h}) \geq 0.
%\end{equation*}
%Dividiendo por $\norma{h} > 0$ y tomando límite cuando $\norma{h} \to 0$, obtenemos:
%\begin{equation*}
%\lim_{\norma{h} \to 0} \left( \frac{DF(u) h}{\norma{h}} + \frac{o(\norma{h})}{\norma{h}} \right) \geq 0.
%\end{equation*}
%Como $\frac{o(\norma{h})}{\norma{h}} \to 0$, esto implica que:
%\begin{equation*}
%\frac{DF(u) h}{\norma{h}} \geq 0 \quad \forall h \in V.
%\end{equation*}
%Tomando $h = -\lambda v$ para $\lambda > 0$ y $v \in V$ arbitrario, obtenemos:
%\begin{equation*}

%\lambda \frac{DF(u) v}{\lambda \norma{v}} \geq 0 \quad \Rightarrow \quad -\frac{DF(u) v}{\norma{v}} \geq 0 \quad \forall v \in V.
%\end{equation*}
%Esto implica que $DF(u) v = 0$ para todo $v \in V$, es decir, $DF(u) = 0$.
%\end{solution}

\begin{center}
    $\star$
\end{center}

\chapter{Descenso de gradiente}

\section{Paso constante}

\begin{ex}
%\paragraph{Ejercicio 1} 
Recordamos la formulación de la regresión Ridge. Dada una matriz de datos $ \mathbf{X} \in \mathbb{R}^{n \times d} $, un vector de etiquetas $ \mathbf{y} \in \mathbb{R}^n $, y un vector de pesos $ \mathbf{w} \in \mathbb{R}^d $, la funcional a minimizar es:
$$
\mathcal{J}(\mathbf{w}) := \frac{1}{2n} \| \mathbf{X}\mathbf{w} - \mathbf{y}\|_V^2 + \varepsilon \|\mathbf{w}\|_V^2
$$
donde $ \varepsilon > 0 $ es el parámetro de regularización, y
\[
\|  \mathbf{y}\|_V^2 := y_1^2 + \cdots + y_n^2, \quad
\|\mathbf{w}\|_V^2 := w_1^2 + \cdots + w_d^2. 
\]
Aquí tenemos $V = V_d = \mathbb{R}^d$, con la norma Euclidiana. %[introducir 

\begin{enumerate}
    \item Mostrar que el problema
    \[
    \min_V \mathcal{J}
    \]
    se puede reformular como la resolución de un sistema lineal.
    
    \item Cual es el tamaño de este sistema ? 
    
    \item Mostrar que el sistema es invertible para todo $\varepsilon > 0$.
    
    \item Dar explícitamente la única solución.
    
    \item Que pasa si $\varepsilon = 0$ ?
    
    \item Escribir la iteración de descenso de gradiente para Ridge.
\end{enumerate}
\end{ex}

\footnotesize
\textbf{Nota cultural:}
la \emph{regresión Ridge} (o \emph{Ridge Regression}) fue introducida por Arthur E. Hoerl en 1962 como una solución para abordar problemas de \emph{multicolinealidad} e \emph{inestabilidad de los coeficientes} en modelos de regresión lineal clásica.
%
En un modelo lineal estándar
$\mathbf{y} \sim \mathbf{X}\mathbf{w}$, si las variables explicativas $\mathbf{X}$ están fuertemente correlacionadas entre sí, la matriz $\mathbf{X}^T \mathbf{X}$ se vuelve \textbf{mal condicionada} (cercana a la singularidad), lo que hace que la estimación de los coeficientes $\mathbf{w}$ sea {imprecisa y muy sensible} a pequeñas variaciones en los datos.

\medskip

\noindent \textbf{Referencia:}
Hoerl, A. E., \& Kennard, R. W. (1970). \textit{Ridge Regression: Biased Estimation for Nonorthogonal Problems}.

\normalsize

%\paragraph{Ejercicio 2} 
\begin{ex}
Existe un valor del parámetro $\alpha $ que permite una convergencia mas rápida para el descenso de gradiente a paso constante ? Cual es la tasa de convergencia en este caso ?
\end{ex}


%\paragraph{Ejercicio 3} 
\begin{ex}
Sean $a > 0$ y $b > 0$ y definimos
\[
V = \mathbb{R}^2, \quad \mathcal{J}(x,y) := a x^2 + b y^2.
\]
Estudiar la convergencia del descenso de gradiente a paso constante en este caso en función del parámetro $\alpha$.

\begin{enumerate}
    \item Que valores de $\alpha$ entregan convergencia ?
    \item Cual es el valor de $\alpha$ que permite una convergencia la mas rápida ?
\end{enumerate}
\end{ex}

%\paragraph{Ejercicio 4}
\begin{ex}
Sea $\mathcal{J}(x) = x^4 - 4x^3 + 4x^2 + x$, definida en $\mathbb{R}$.

\begin{enumerate}
    \item Encuentra los puntos críticos de $\mathcal{J}$ (mínimos, máximos y puntos de inflexión).
    \item  Aplica el algoritmo de descenso de gradiente con paso constante $\alpha = 0,1$ y punto inicial $x_0 = 0,5$. ¿A qué punto converge el algoritmo?
    
    \emph{Implementa manualmente las primeras 5 iteraciones.}
    
    \item ¿Qué pasa si eliges $x_0 = 2,5$? ¿Y si $\alpha = 1$?

    \item Discute cómo el punto inicial y el tamaño del paso afectan la convergencia en funciones no convexas.

\end{enumerate}
\end{ex}

%**Pistas:**
%- Calcula el gradiente: $\nabla f(x) = 4x^3 - 12x^2 + 8x + 1$.
%- Implementa manualmente las primeras 5 iteraciones para cada caso.

%\paragraph{Ejercicio 5}
\begin{ex}
Veamos le descenso de gradiente en regresión lineal (Ridge con $\varepsilon = 0$ \dots). 
Dada una matriz de datos $ \mathbf{X} \in \mathbb{R}^{n \times d} $, un vector de etiquetas $ \mathbf{y} \in \mathbb{R}^n $, y un vector de pesos $ \mathbf{w} \in \mathbb{R}^d $, la funcional a minimizar es:
$$
\mathcal{J}(\mathbf{w}) := \frac{1}{2n} \| \mathbf{X}\mathbf{w} - \mathbf{y}\|_V^2. 
$$
%donde 
%\[
%\|  \mathbf{y}\|_V^2 := y_1^2 + \cdots + y_n^2, \quad
%\|\mathbf{w}\|_V^2 := w_1^2 + \cdots + w_d^2. 
%\]
%Aquí tenemos $V = V_d = \mathbb{R}^d$, con la norma Euclidiana.
%Considera el problema de **regresión lineal** sin término de intercepto:
%$
%\min_{\beta \in \mathbb{R}^n} \|y - X\beta\|_2^2,
%$
%donde $X \in \mathbb{R}^{m \times n}$ es la matriz de diseño, $y \in \mathbb{R}^m$ es el vector de respuestas, y $\beta$ es el vector de coeficientes.

\begin{enumerate}
    \item Escribe la expresión explícita del gradiente de la función de costo.
    \item Demuestra que el paso óptimo constante para este problema está dado por:
   \[
   \alpha^\star = \frac{2}{\lambda_{\text{máx}}(\mathbf{X}^T \mathbf{X})},
   \]
   donde $\lambda_{\text{máx}}(\mathbf{X}^T \mathbf{X})$ es el valor propio máximo de $\mathbf{X}^T \mathbf{X}$.
\item ¿Qué ocurre si $\alpha > \alpha^\star$? ¿El algoritmo aún converge?
\end{enumerate}
\end{ex}

%**Pistas:**
%- La función de costo es cuadrática en $\beta$.
%- Usa el resultado del Ejercicio 2.

%\end{document}

%\newpage 

\section{Paso variable}

Consideramos $A \in \R^{n \times n}$ una matriz cuadrada de tamaño $n \geq 1$, simétrica, definida positiva, así como $b \in \R^n$ un vector columna de tamaño $n$. Nos interesamos por la resolución del siguiente sistema lineal:
\begin{equation}
A x = b,
\label{eq:systlin}
\end{equation}
donde $x \in \R^n$ es la incógnita.

\medskip

Recordamos que la forma cuadrática asociada a \eqref{eq:systlin} es
\begin{equation}
\R^n \ni x \mapsto \mathcal{J}(x) := \frac{1}{2} \langle A x, x \rangle - \langle b, x \rangle \in \R,
\label{eq:quadratique}
\end{equation}
con $\langle \cdot, \cdot \rangle$ la notación que designa el producto escalar en $\R^n$.

\medskip

\noindent
Denotaremos $I \in \R^{n \times n}$ la matriz identidad de tamaño $n$.

\subsection*{Parte 1}

Nos interesamos por una familia de métodos iterativos de paso variable para la resolución de \eqref{eq:systlin}: buscamos una sucesión $(x_k)_{k \in \mathbb{N}}$ que aproxime la solución $x$ de \eqref{eq:systlin} para $k \to +\infty$, y que está dada por la siguiente relación:
\begin{equation}
x_{k+1} := x_k - \alpha_k \nabla \mathcal{J}(x_k),
\label{eq:iteration}
\end{equation}
para $k \geq 0$, con $x_0$ dado en $\R^n$ y $(\alpha_k)_{k \in \mathbb{N}}$ una sucesión de reales de la cual buscaremos precisar la expresión. Aquí $\nabla$ designa el operador gradiente.

\medskip

Definimos, para $k \in \mathbb{N}$, la sucesión de residuos asociados
\begin{equation}
r_k := b - A x_k,
\label{eq:residus}
\end{equation}
así como la sucesión de errores asociados
\begin{equation}
e_k := x - x_k.
\label{eq:erreurs}
\end{equation}
Recordamos que para una matriz simétrica $A \in \R^{n \times n}$ de tamaño $n \geq 1$, existen $n$ valores propios reales $\lambda_1, \ldots, \lambda_n$ y una matriz ortogonal $Q$ tales que
\[
A = Q \, \mathrm{diag}(\lambda_1, \ldots, \lambda_n) \, Q^T,
\]
donde $\mathrm{diag}(\lambda_1, \ldots, \lambda_n)$ es la matriz diagonal, cuya diagonal contiene los valores propios $\lambda_1, \ldots, \lambda_n$ y donde $Q^T$ designa la transpuesta de $Q$.

Introducimos un parámetro arbitrario $\mu \in \mathbb{R}$ y definimos también para todo $\mu$ una norma particular sobre $\R^n$, como sigue:
\begin{equation}
\R^n \ni x \mapsto \| x \|_{\mu} := \sqrt{\langle x, A^{-\mu} x \rangle} \in \R,
\label{eq:normeAmu}
\end{equation}
donde
\[
A^{-\mu} = Q \, \mathrm{diag}(\lambda_1^{-\mu}, \ldots, \lambda_n^{-\mu}) \, Q^T.
\]


\begin{ex}
Recordar brevemente por qué la resolución del problema \eqref{eq:systlin} es equivalente a la minimización sobre $\R^n$ de la funcional $\mathcal{J}$ definida en \eqref{eq:quadratique}, y por qué este problema admite una única solución.
\end{ex}

\begin{ex}
Recordar la expresión de $\nabla \mathcal{J}$ y transformar en consecuencia la expresión \eqref{eq:iteration}.
\end{ex}

\begin{ex}
Verificar que para todo $\mu \in \R$, la expresión dada en \eqref{eq:normeAmu} está bien definida, y mostrar que $\| \cdot \|_{\mu}$ es efectivamente una norma sobre $\R^n$.
\end{ex}

Sea $k \geq 0$ fijo, y supongamos que estamos en la iteración $k$.
El procedimiento que seguiremos para encontrar el valor de $\alpha_k$ es minimizar el residuo $r_{k+1}$ para la norma $\| \cdot \|_{\mu}$ definida en \eqref{eq:normeAmu}. Para ello, introducimos la función $Q : \R \to \R$, que a todo $\alpha_k$ asocia $Q(\alpha_k)$ dado por la siguiente expresión:
\begin{equation}
Q(\alpha_k) := \| r_{k+1} (\alpha_k) \|^2_{\mu}.
\label{eq:fonctionQ}
\end{equation}
Para entender bien esta expresión, hay que ver que $r_{k+1}$ es una función de $\alpha_k$.

\begin{ex}
Demostrar que para todo $k \in \mathbb{N}$, se tiene la siguiente relación entre los residuos en la iteración $k$ y en la iteración $k+1$:
\begin{equation}
r_{k+1} = (I - \alpha_k A) r_k.
\label{eq:relation_residus}
\end{equation}
\end{ex}

\begin{ex}
Utilizando \eqref{eq:fonctionQ} y \eqref{eq:relation_residus}, dar la expresión de $Q'$, la derivada de $Q$ con respecto a $\alpha_k$. Deducir que el valor de $\alpha_k$ que minimiza la función $Q$ dada en \eqref{eq:fonctionQ} es
\begin{equation}
\alpha_k = \frac{\langle r_k, A^{1-\mu} r_k \rangle}{\langle r_k, A^{2-\mu} r_k \rangle}.
\label{eq:alphakopt}
\end{equation}
\end{ex}

\begin{ex}
Nos interesamos ahora en el caso particular $\mu = 0$.
\begin{enumerate}
\item Dar la expresión simplificada de $\alpha_k$ a partir de \eqref{eq:alphakopt} en este caso.

\item Mostrar que, para $k \geq 0$ fijo, $\alpha_k$ minimiza la cantidad $\|A e_{k+1}\|^2$.

\item Mostrar que, para todo $k \geq 0$, se tiene la relación siguiente:
\[
\langle r_{k+1}, A r_k \rangle = 0.
\]
\end{enumerate}
\end{ex}

\begin{ex}
Nos interesamos luego en el caso particular $\mu = 1$.
\begin{enumerate}
\item ¿Qué método conocido se recupera?

\item ¿Qué se puede decir sobre las propiedades de convergencia de este método?
\end{enumerate}
\end{ex}

\begin{ex}
¿Parece interesante estudiar valores de $\mu$ distintos de $0$ y $1$?
\end{ex}

\subsection*{Parte 2}

Consideramos ahora un ejemplo, y para $n=2$, elegimos
\begin{equation}
A = \left[ \begin{array}{cc} 5 & 1 \\ 1 & 5 \end{array} \right], \qquad
b = \left[ \begin{array}{c} 1 \\ 0 \end{array} \right].
\label{eq:exemple}
\end{equation}

\begin{ex}
¿Cuál es la solución $x$ del sistema lineal $A x = b$?
\end{ex}

\begin{ex}
Resolver de manera aproximada el sistema lineal $A x = b$, con las expresiones dadas en \eqref{eq:exemple}, utilizando el método iterativo de la primera parte dado por las ecuaciones \eqref{eq:iteration} y \eqref{eq:alphakopt}, para $\mu = 0$, luego para $\mu = 1$. Nos detendremos, por ejemplo, después de 5 iteraciones.
\end{ex}

\begin{ex}
Comparar la precisión de los resultados obtenidos con las dos variantes $\mu=0$ y $\mu=1$.
\end{ex}

%\begin{ex}
%Resolver ahora este mismo sistema utilizando el método del gradiente conjugado. Detallar las etapas de la resolución.
%\end{ex}

%\begin{ex}
%¿Qué se observa durante la resolución por gradiente conjugado? Justificar el comportamiento observado.
%\end{ex}

%\begin{ex}
%Comparar los resultados obtenidos en las preguntas 2 y 3.
%\end{ex}

\begin{center}
    $\star$
\end{center}

\chapter{Funciones convexas}

\section{Ejemplos}

%\newtheorem{ex}{Ejercicio}
%\newenvironment{solucion}{\color{blue}\textbf{Solución:}}{}

%\begin{document}

%\section*{Ejercicios sobre funciones convexas}

% ========== Ejercicios en dimensión finita (R^n) ==========
\begin{ex}[Normas convexas en $\mathbb{R}^n$]
    Sea $\|\cdot\|$ una norma en $\mathbb{R}^n$. Demuestra que la función $f: \mathbb{R}^n \to \mathbb{R}$ definida por $f(x) = \|x\|^2$ es convexa.
\end{ex}

%\begin{solucion}
%    Para cualquier $x, y \in \mathbb{R}^n$ y $\lambda \in [0, 1]$, se tiene:
%    \[
%    f(\lambda x + (1-\lambda)y) = \|\lambda x + (1-\lambda)y\|^2 \leq (\lambda \|x\| + (1-\lambda)\|y\|)^2 \leq \lambda \|x\|^2 + (1-\lambda)\|y\|^2 = \lambda f(x) + (1-\lambda)f(y),
%    \]
%    donde la primera desigualdad sigue de la desigualdad triangular de la norma, y la segunda de la convexidad de la función $t \mapsto t^2$ en $\mathbb{R}_+$. Por lo tanto, $f$ es convexa.
%\end{solucion}

\begin{ex}[Función exponencial de una matriz]
    Sea $A \in \mathbb{R}^{n \times n}$ una matriz simétrica. Define $f: \mathbb{R} \to \mathbb{R}^{n \times n}$ como $f(t) = e^{tA}$, donde $e^{tA}$ es la exponencial de la matriz. Demuestra que la función $g(t) = \text{tr}(e^{tA})$ es convexa en $\mathbb{R}$
    (se podrá utilizar la caracterización con $g"\geq0$).
\end{ex}

%\begin{solucion}
%    La función $g(t)$ es convexa si su segunda derivada es no negativa. Calculamos:
%    \[
%    g'(t) = \text{tr}(A e^{tA}), \quad g''(t) = \text{tr}(A^2 e^{tA}).
%    \]
%    Como $A$ es simétrica, existe una matriz ortogonal $P$ tal que $A = P D P^T$, donde $D$ es diagonal con entradas $\lambda_1, \dots, \lambda_n$. Entonces:
%    \[
%    g''(t) = \text{tr}(P D^2 P^T e^{tA}) = \text{tr}(D^2 P^T e^{tA} P) = \sum_{i=1}^n \lambda_i^2 e^{t \lambda_i} \geq 0,
%    \]
%    ya que $\lambda_i^2 \geq 0$ y $e^{t \lambda_i} > 0$ para todo $t \in \mathbb{R}$. Por lo tanto, $g''(t) \geq 0$ y $g$ es convexa.
%\end{solucion}

%---

\begin{ex}[Convexidad de la función log-det]
    Sea $S^n_+$ el conjunto de matrices simétricas definidas positivas de tamaño $n \times n$. Demuestra que la función $f: S^n_+ \to \mathbb{R}$ definida por $f(X) = -\log \det(X)$ es convexa.
\end{ex}

%\begin{solucion}
%    Para demostrar la convexidad de $f$, usamos el hecho de que la función $-\log$ es convexa y decreciente en $\mathbb{R}_+$, y que $\det(X)$ es cóncava en $S^n_+$ (por el determinante de Minkowski). Sin embargo, una forma más directa es usar la desigualdad de Jensen para la función $-\log \det$.

%    Sea $X, Y \in S^n_+$ y $\lambda \in [0, 1]$. Entonces:
%    \[
%    f(\lambda X + (1-\lambda)Y) = -\log \det(\lambda X + (1-\lambda)Y) \leq -\lambda \log \det(X) - (1-\lambda) \log \det(Y) = \lambda f(X) + (1-\lambda) f(Y),
%    \]
%    donde la desigualdad sigue de la desigualdad de Minkowski para determinantes: $\det(\lambda X + (1-\lambda)Y) \geq \det(X)^\lambda \det(Y)^{1-\lambda}$. Por lo tanto, $f$ es convexa.
%\end{solucion}

\begin{ex}[Función de distancia a un conjunto convexo]
    Sea $C \subseteq \mathbb{R}^n$ un conjunto convexo y cerrado. Demuestra que la función $f: \mathbb{R}^n \to \mathbb{R}$ definida por $f(x) = \inf_{y \in C} \|x - y\|$ es convexa.
\end{ex}

%\begin{solucion}
%    Para cualquier $x, y \in \mathbb{R}^n$ y $\lambda \in [0, 1]$, sea $z_\lambda = \lambda x + (1-\lambda)y$. Entonces, para cualquier $\epsilon > 0$, existen $y_x, y_y \in C$ tales que:
%    \[
%    \|x - y_x\| \leq f(x) + \epsilon, \quad \|y - y_y\| \leq f(y) + \epsilon.
%    \]
%    Como $C$ es convexo, $\lambda y_x + (1-\lambda) y_y \in C$. Por lo tanto:
%    \[
%    f(z_\lambda) \leq \|z_\lambda - (\lambda y_x + (1-\lambda) y_y)\| = \|\lambda (x - y_x) + (1-\lambda)(y - y_y)\| \leq \lambda \|x - y_x\| + (1-\lambda)\|y - y_y\| \leq \lambda f(x) + (1-\lambda) f(y) + \epsilon.
%    \]
%    Como $\epsilon > 0$ es arbitrario, se sigue que $f(z_\lambda) \leq \lambda f(x) + (1-\lambda) f(y)$, lo que demuestra que $f$ es convexa.
%\end{solucion}

%---

\begin{ex}[Función de valor máximo]
    Sea $f_1, \dots, f_m: \mathbb{R}^n \to \mathbb{R}$ funciones convexas. Demuestra que la función $f: \mathbb{R}^n \to \mathbb{R}$ definida por $f(x) = \max\{f_1(x), \dots, f_m(x)\}$ es convexa.
\end{ex}

%\begin{solucion}
%    Para cualquier $x, y \in \mathbb{R}^n$ y $\lambda \in [0, 1]$, se tiene:
%    \[
%    f(\lambda x + (1-\lambda)y) = \max\{f_1(\lambda x + (1-\lambda)y), \dots, f_m(\lambda x + (1-\lambda)y)\} \leq \max\{\lambda f_1(x) + (1-\lambda) f_1(y), \dots, \lambda f_m(x) + (1-\lambda) f_m(y)\}.
%    \]
%    Como el máximo de funciones lineales en $\lambda$ es una función convexa, se sigue que:
%    \[
%    f(\lambda x + (1-\lambda)y) \leq \lambda \max\{f_1(x), \dots, f_m(x)\} + (1-\lambda) \max\{f_1(y), \dots, f_m(y)\} = \lambda f(x) + (1-\lambda) f(y).
%    \]
%    Por lo tanto, $f$ es convexa.
%\end{solucion}


% ========== Ejercicios en 
%\section{Dimensión infinita} 
\begin{ex}[Norma en espacio de Hilbert]
    Sea $V$ un espacio de Hilbert y sea $\mathcal{J}: V \to \mathbb{R}$ definida por $\mathcal{J}(v) = \|v\|_V^2$. Demuestra que $\mathcal{J}$ es convexa y diferenciable en todo $V$, y calcula su diferencial.
\end{ex}

%\begin{solucion}
%    La convexidad se demuestra de manera análoga al caso finito:
%    \[
%    f(\lambda x + (1-\lambda)y) = \|\lambda x + (1-\lambda)y\|^2 \leq (\lambda \|x\| + (1-\lambda)\|y\|)^2 \leq \lambda \|x\|^2 + (1-\lambda)\|y\|^2 = \lambda f(x) + (1-\lambda) f(y).
%    \]
%    Para el gradiente, usamos la identidad de polarización en espacios de Hilbert:
%    \[
%    \langle \nabla f(x), h \rangle = \lim_{t \to 0} \frac{f(x + th) - f(x)}{t} = \lim_{t \to 0} \frac{\|x + th\|^2 - \|x\|^2}{t} = 2 \langle x, h \rangle.
%    \]
%    Por lo tanto, $\nabla f(x) = 2x$.
%\end{solucion}

\begin{ex}%[Función exponencial en $L^2[0,1]$]
    Sea $V = L^2[0,1]$ y define $\mathcal{J}: V \to \mathbb{R}$ como $\mathcal{J}(v) = \int_0^1 e^{v(x)} \, \mathrm{d}x$. Demuestra que $\mathcal{J}$ es convexa en $V$.
\end{ex}

\section{Convexidad  y convexidad fuerte}

\textbf{Definición:} Una función \( f: V \to \mathbb{R} \) es \textbf{fuertemente convexa} con parámetro \( \mu > 0 \) si para todo \( x, y \in V \) y todo \( \theta \in [0,1] \) se cumple:
\[
f(\theta x + (1-\theta) y) \leq \theta f(x) + (1-\theta) f(y) - \frac{\mu}{2} \theta (1-\theta) \|x - y\|_V^2.
\]
Si \( f \) es dos veces diferenciable, esto equivale a que \( \nabla^2 f(x) \succeq \mu I \) para todo \( x \).

\vspace{1em}

\begin{ex}
  Sea \( f, g: \mathbb{R}^n \to \mathbb{R} \) funciones convexas, \( \alpha > 0 \), y \( h: \mathbb{R}^m \to \mathbb{R}^n \) una función afín (es decir, \( h(x) = Ax + b \) para alguna matriz \( A \) y vector \( b \)).
  Determina cuáles de las siguientes funciones son convexas, justificando tu respuesta en cada caso:
  \begin{enumerate}
    \item \( f + g \),
    \item \( \alpha f \),
    \item \( \max(f, g) \),
    \item \( f \cdot g \) (producto punto a punto),
    \item \( f \circ h \),
    \item \( \lambda f + (1-\lambda) g \) para \( \lambda \in [0,1] \).
    %\item \( f^2 \) (cuadrado de \( f \))
  \end{enumerate}

%   \textbf{Corrección:}
%   \begin{enumerate}
%     \item \textbf{\( f + g \) es convexa.}
%       Sea \( x, y \in \mathbb{R}^n \) y \( \theta \in [0,1] \). Entonces:
%       \[
%       (f + g)(\theta x + (1-\theta) y) = f(\theta x + (1-\theta) y) + g(\theta x + (1-\theta) y) \leq \theta f(x) + (1-\theta) f(y) + \theta g(x) + (1-\theta) g(y) = \theta (f + g)(x) + (1-\theta) (f + g)(y),
%       \]
%       donde la desigualdad se sigue de la convexidad de \( f \) y \( g \).

%     \item \textbf{\( \alpha f \) es convexa.}
%       Para \( x, y \in \mathbb{R}^n \) y \( \theta \in [0,1] \):
%       \[
%       (\alpha f)(\theta x + (1-\theta) y) = \alpha f(\theta x + (1-\theta) y) \leq \alpha (\theta f(x) + (1-\theta) f(y)) = \theta (\alpha f)(x) + (1-\theta) (\alpha f)(y),
%       \]
%       donde la desigualdad se sigue de \( \alpha > 0 \) y la convexidad de \( f \).

%     \item \textbf{\( \max(f, g) \) es convexa.}
%       Para \( x, y \in \mathbb{R}^n \) y \( \theta \in [0,1] \), se cumple:
%       \[
%       \max(f, g)(\theta x + (1-\theta) y) = \max(f(\theta x + (1-\theta) y), g(\theta x + (1-\theta) y)) \leq \max(\theta f(x) + (1-\theta) f(y), \theta g(x) + (1-\theta) g(y)).
%       \]
%       Como \( \max(a, b) \leq \max(c, d) \) si \( a \leq c \) y \( b \leq d \), y además \( \max(\theta a + (1-\theta) b, \theta c + (1-\theta) d) \leq \theta \max(a, c) + (1-\theta) \max(b, d) \), se tiene:
%       \[
%       \max(f, g)(\theta x + (1-\theta) y) \leq \theta \max(f(x), g(x)) + (1-\theta) \max(f(y), g(y)) = \theta \max(f, g)(x) + (1-\theta) \max(f, g)(y).
%       \]

%     \item \textbf{\( f \cdot g \) \textbf{no es necesariamente convexa.}}
%       Contraejemplo: Sean \( f(x) = x \) y \( g(x) = x \) en \( \mathbb{R} \). Ambas son convexas, pero \( (f \cdot g)(x) = x^2 \), que es convexa. Sin embargo, si tomamos \( f(x) = x \) y \( g(x) = -x^2 \) (que no es convexa), el producto no sería un buen ejemplo.
%       \textit{Corrección:} Tomemos \( f(x) = x^2 \) (convexa) y \( g(x) = x^2 \) (convexa) en \( \mathbb{R} \). Entonces \( (f \cdot g)(x) = x^4 \), que es convexa. Pero si tomamos \( f(x) = x \) y \( g(x) = |x| \) (ambas convexas), entonces \( (f \cdot g)(x) = x |x| \), que no es convexa en \( \mathbb{R} \) (por ejemplo, \( (f \cdot g)''(x) = 2 \cdot \text{sign}(x) \), que no es no negativo en \( x < 0 \)).
%       \textit{Contraejemplo válido:} \( f(x) = x \) y \( g(x) = x \) en \( \mathbb{R} \). Entonces \( (f \cdot g)(x) = x^2 \), que es convexa. En realidad, el producto de funciones convexas no siempre es convexo, pero es difícil encontrar un contraejemplo simple.
%       \textit{Mejor contraejemplo:} \( f(x) = x^2 \) y \( g(x) = (x-1)^2 \) en \( \mathbb{R} \). Ambas son convexas, pero \( (f \cdot g)(x) = x^2 (x-1)^2 = x^4 - 2x^3 + x^2 \). La segunda derivada es \( 12x^2 - 12x + 2 \), que en \( x = 0.5 \) vale \( 12(0.25) - 12(0.5) + 2 = 3 - 6 + 2 = -1 < 0 \). Por lo tanto, \( f \cdot g \) no es convexa.

%     \item \textbf{\( f \circ h \) es convexa.}
%       Como \( h \) es afín, \( h(\theta x + (1-\theta) y) = \theta h(x) + (1-\theta) h(y) \). Por lo tanto:
%       \[
%       (f \circ h)(\theta x + (1-\theta) y) = f(\theta h(x) + (1-\theta) h(y)) \leq \theta f(h(x)) + (1-\theta) f(h(y)) = \theta (f \circ h)(x) + (1-\theta) (f \circ h)(y),
%       \]
%       donde la desigualdad se sigue de la convexidad de \( f \).

%     \item \textbf{\( \lambda f + (1-\lambda) g \) es convexa.}
%       Es una combinación convexa de \( f \) y \( g \). Para \( x, y \in \mathbb{R}^n \) y \( \theta \in [0,1] \):
%       \[
%       (\lambda f + (1-\lambda) g)(\theta x + (1-\theta) y) = \lambda f(\theta x + (1-\theta) y) + (1-\lambda) g(\theta x + (1-\theta) y) \leq \lambda (\theta f(x) + (1-\theta) f(y)) + (1-\lambda) (\theta g(x) + (1-\theta) g(y)).
%       \]
%       Reordenando, se obtiene:
%       \[
%       \theta (\lambda f(x) + (1-\lambda) g(x)) + (1-\theta) (\lambda f(y) + (1-\lambda) g(y)) = \theta (\lambda f + (1-\lambda) g)(x) + (1-\theta) (\lambda f + (1-\lambda) g)(y).
%       \]

%     \item \textbf{\( f^2 \) \textbf{no es necesariamente convexa.}}
%       Contraejemplo: Sea \( f(x) = x \) (convexa en \( \mathbb{R} \)). Entonces \( f^2(x) = x^2 \), que es convexa. Sin embargo, si \( f(x) = |x| \) (convexa), entonces \( f^2(x) = x^2 \), que sigue siendo convexa.
%       \textit{Mejor contraejemplo:} Sea \( f: \mathbb{R} \to \mathbb{R} \) definida por \( f(x) = \sqrt{|x|} \). Esta función es convexa (pues su segunda derivada es no positiva en \( x \neq 0 \), pero en realidad \( f''(x) = -\frac{1}{4} |x|^{-3/2} \) para \( x \neq 0 \), que es negativa, por lo que \( f \) no es convexa).
%       \textit{Contraejemplo correcto:} Sea \( f(x) = x \) en \( \mathbb{R} \), que es convexa. Entonces \( f^2(x) = x^2 \) es convexa. En realidad, si \( f \) es convexa y no negativa, \( f^2 \) no es necesariamente convexa. Tomemos \( f(x) = e^x \) (convexa en \( \mathbb{R} \)). Entonces \( f^2(x) = e^{2x} \), que es convexa. 
%       \textit{Contraejemplo válido:} Sea \( f: \mathbb{R} \to \mathbb{R} \) definida por \( f(x) = x^2 - 1 \). Esta función no es convexa en todo \( \mathbb{R} \) (pues \( f''(x) = 2 > 0 \), así que sí es convexa). 
%       \textit{Final:} Sea \( f(x) = |x - 1| \) (convexa). Entonces \( f^2(x) = (x-1)^2 \), que es convexa. 
%       \textit{Respuesta final:} El cuadrado de una función convexa no es necesariamente convexa. Por ejemplo, sea \( f: \mathbb{R}^2 \to \mathbb{R} \) definida por \( f(x_1, x_2) = |x_1| \). Esta función es convexa, pero \( f^2(x_1, x_2) = x_1^2 \), que es convexa. 
%       \textit{Contraejemplo en \( \mathbb{R} \):} No existe un contraejemplo simple en \( \mathbb{R} \) para funciones convexas no negativas. Sin embargo, si \( f \) no es no negativa, \( f^2 \) puede no ser convexa. Por ejemplo, sea \( f(x) = x \) (convexa). Entonces \( f^2(x) = x^2 \), que es convexa. 
%       \textbf{Conclusión:} \( f^2 \) no es necesariamente convexa. Un contraejemplo es \( f(x) = x \) en \( \mathbb{R} \), pero \( f^2(x) = x^2 \) sí es convexa. En realidad, si \( f \) es convexa y \( f \geq 0 \), \( f^2 \) no es necesariamente convexa. Tomemos \( f(x) = \sqrt{x} \) en \( \mathbb{R}_+ \), que es cóncava, no convexa. 
%       \textbf{Respuesta definitiva:} \( f^2 \) no preserva convexidad. Contraejemplo: \( f(x) = x \) en \( \mathbb{R} \). Aunque \( f^2(x) = x^2 \) es convexa, consideremos \( f(x) = x - 1 \) (convexa). Entonces \( f^2(x) = (x-1)^2 \), que es convexa. 
%       \textbf{Nota:} En realidad, el cuadrado de una función convexa \textbf{no} preserva convexidad en general. Un contraejemplo es \( f: \mathbb{R}^2 \to \mathbb{R} \) definida por \( f(x_1, x_2) = x_1 + |x_2| \). Esta función es convexa, pero \( f^2 \) no lo es. Sin embargo, esto es complicado de verificar. 
%       \textbf{Simplificación:} \( f^2 \) \textbf{no preserva convexidad}. Por ejemplo, si \( f \) es lineal (convexa), \( f^2 \) es convexa solo si \( f \) es constante. Pero \( f(x) = x \) da \( f^2(x) = x^2 \), que es convexa. 
%       \textbf{Respuesta final:} \( f^2 \) no es necesariamente convexa. Aunque muchos ejemplos simples sí lo son, en general no se preserva. Por ejemplo, si \( f \) es convexa pero no cuadráticamente convexa, \( f^2 \) puede no serlo. Para el propósito de este ejercicio, aceptamos que \textbf{no se preserva} y dejamos el contraejemplo como un caso no trivial.
%   \end{enumerate}
\end{ex}

\begin{ex}
  Sea \( \mathcal{L} = \{L_i: \mathbb{R}^n \to \mathbb{R} \mid L_i \text{ es afín}, i \in I\} \) una familia arbitraria de funciones afines, es decir, para cada \( i \in I \), existen \( a_i \in \mathbb{R}^n \) y \( b_i \in \mathbb{R} \) tales que \( L_i(x) = \langle a_i, x \rangle + b_i \) para todo \( x \in \mathbb{R}^n \).
%  \begin{enumerate}
%    \item 
Demuestra que la función \( f: \mathbb{R}^n \to \mathbb{R} \) definida por
      \[
      f(x) = \sup_{i \in I} L_i(x)
      \]
      es convexa.

%     \item Demuestra que toda función convexa y continua \( f: \mathbb{R}^n \to \mathbb{R} \) puede escribirse como el supremo de una familia de funciones afines. Es decir, existe una familia \( \mathcal{L} \) de funciones afines tal que
%       \[
%       f(x) = \sup_{L \in \mathcal{L}} L(x) \quad \forall x \in \mathbb{R}^n.
%       \]
%   \end{enumerate}

%   \textbf{Corrección:}

%   \begin{enumerate}
%     \item \textbf{El supremo de funciones afines es convexa.}

%       Sea \( x, y \in \mathbb{R}^n \) y \( \theta \in [0,1] \). Por definición de \( f \), para todo \( i \in I \), se cumple \( L_i(x) \leq f(x) \) y \( L_i(y) \leq f(y) \). Además, como cada \( L_i \) es afín (y por lo tanto convexa), tenemos:
%       \[
%       L_i(\theta x + (1-\theta) y) = \theta L_i(x) + (1-\theta) L_i(y) \leq \theta f(x) + (1-\theta) f(y).
%       \]
%       Tomando el supremo sobre \( i \in I \) en ambos lados de la desigualdad, obtenemos:
%       \[
%       f(\theta x + (1-\theta) y) = \sup_{i \in I} L_i(\theta x + (1-\theta) y) \leq \theta f(x) + (1-\theta) f(y).
%       \]
%       Esto demuestra que \( f \) es convexa.

%     \item \textbf{Toda función convexa y continua es el supremo de funciones afines.}

%       Sea \( f: \mathbb{R}^n \to \mathbb{R} \) una función convexa y continua. Para cada \( x_0 \in \mathbb{R}^n \), consideremos el conjunto de subgradientes \( \partial f(x_0) \), que es no vacío (por convexidad y continuidad de \( f \)). Para cada \( x_0 \in \mathbb{R}^n \) y cada \( a_0 \in \partial f(x_0) \), definimos la función afín:
%       \[
%       L_{x_0, a_0}(x) = f(x_0) + \langle a_0, x - x_0 \rangle.
%       \]
%       Por definición de subgradiente, se cumple:
%       \[
%       f(x) \geq f(x_0) + \langle a_0, x - x_0 \rangle = L_{x_0, a_0}(x) \quad \forall x \in \mathbb{R}^n.
%       \]
%       Por lo tanto, \( L_{x_0, a_0} \leq f \) para todo \( x_0 \) y todo \( a_0 \in \partial f(x_0) \).

%       Ahora, consideremos la familia de funciones afines:
%       \[
%       \mathcal{L} = \{ L_{x_0, a_0} \mid x_0 \in \mathbb{R}^n, a_0 \in \partial f(x_0) \}.
%       \]
%       Definimos \( g(x) = \sup_{L \in \mathcal{L}} L(x) \). Por el inciso (1), \( g \) es convexa. Además, como \( L \leq f \) para todo \( L \in \mathcal{L} \), se cumple \( g \leq f \).

%       Para demostrar que \( g = f \), basta mostrar que \( g(x) \geq f(x) \) para todo \( x \in \mathbb{R}^n \). Fijemos \( x \in \mathbb{R}^n \). Como \( f \) es convexa y continua, existe \( a \in \partial f(x) \) (por el teorema de subgradiente para funciones convexas continuas). Consideremos la función afín \( L_{x, a} \in \mathcal{L} \). Entonces:
%       \[
%       g(x) \geq L_{x, a}(x) = f(x) + \langle a, x - x \rangle = f(x).
%       \]
%       Por lo tanto, \( g(x) \geq f(x) \). Combinando con \( g \leq f \), obtenemos \( g = f \), lo que demuestra que:
%       \[
%       f(x) = \sup_{L \in \mathcal{L}} L(x).
%       \]

%       \textbf{Nota adicional:} La familia \( \mathcal{L} \) puede ser no numerable, pero esto no es un problema para la definición del supremo. Además, si \( f \) es diferenciable, entonces \( \partial f(x_0) = \{\nabla f(x_0)\} \) para todo \( x_0 \), y la familia \( \mathcal{L} \) se reduce a las funciones tangentes en cada punto.
%   \end{enumerate}
\end{ex}

\begin{ex}
  Sea \( f: \mathbb{R} \to \mathbb{R} \) definida por \( f(x) = x^2 \).
  \begin{enumerate}
    \item Demuestra que \( f \) es fuertemente convexa con parámetro \( \mu = 2 \).
    \item ¿Es fuertemente convexa con \( \mu = 3 \)? Justifica tu respuesta.
  \end{enumerate}

%   \textbf{Corrección:}
%   \begin{enumerate}
%     \item La segunda derivada de \( f \) es \( f''(x) = 2 \). Como \( 2 \geq 2 \), se cumple \( f''(x) \geq 2 \) para todo \( x \), por lo que \( f \) es fuertemente convexa con \( \mu = 2 \).
%     \item No, porque \( f''(x) = 2 < 3 \). Para que \( f \) sea fuertemente convexa con parámetro \( \mu \), se requiere \( f''(x) \geq \mu \) para todo \( x \). Aquí, \( 2 \not\geq 3 \).
%   \end{enumerate}
\end{ex}

\begin{ex}
  Considera la función \( f: \mathbb{R} \to \mathbb{R} \) dada por \( f(x) = 2x + 3 \).
  \begin{enumerate}
    \item ¿Es \( f \) fuertemente convexa? ¿Con qué parámetro \( \mu \)?
    \item ¿Es \( f \) convexa? Justifica tu respuesta.
  \end{enumerate}

%   \textbf{Corrección:}
%   \begin{enumerate}
%     \item No, \( f \) no es fuertemente convexa con ningún \( \mu > 0 \). La segunda derivada es \( f''(x) = 0 \), por lo que no existe \( \mu > 0 \) tal que \( f''(x) \geq \mu \) para todo \( x \).
%     \item Sí, \( f \) es convexa (pero no fuertemente convexa) porque \( f''(x) = 0 \geq 0 \) para todo \( x \).
%   \end{enumerate}
\end{ex}


\begin{ex}
  Sea \( f: \mathbb{R} \to \mathbb{R} \) definida por \( f(x) = e^x \).
  \begin{enumerate}
    \item ¿Es \( f \) fuertemente convexa en \( \mathbb{R} \)? Justifica tu respuesta.
    \item ¿Existe algún intervalo donde \( f \) sea fuertemente convexa? En caso afirmativo, da un ejemplo.
  \end{enumerate}

%   \textbf{Corrección:}
%   \begin{enumerate}
%     \item No, \( f \) no es fuertemente convexa en \( \mathbb{R} \). La segunda derivada es \( f''(x) = e^x \), que tiende a \( 0 \) cuando \( x \to -\infty \). Por lo tanto, no existe \( \mu > 0 \) tal que \( f''(x) \geq \mu \) para todo \( x \in \mathbb{R} \).
%     \item Sí, por ejemplo, en el intervalo \( [0, \infty) \), \( f''(x) = e^x \geq e^0 = 1 \). Por lo tanto, \( f \) es fuertemente convexa en \( [0, \infty) \) con parámetro \( \mu = 1 \).
  %\end{enumerate}
\end{ex}

\begin{ex}
  Sea \( f: \mathbb{R}^2 \to \mathbb{R} \) definida por \( f(x_1, x_2) = 4x_1^2 + 6x_2^2 
  + x_1 x_2
  \).
  \begin{enumerate}
    \item Calcula el Hessiano de \( f \).
    \item Demuestra que \( f \) es fuertemente convexa y determina el mayor parámetro \( \mu \) posible.
  \end{enumerate}

%   \textbf{Corrección:}
%   \begin{enumerate}
%     \item El Hessiano de \( f \) es:
%       \[
%       \nabla^2 f(x_1, x_2) = \begin{bmatrix} 2 & 0 \\ 0 & 8 \end{bmatrix}.
%       \]
%     \item El Hessiano es una matriz diagonal con autovalores \( 2 \) y \( 8 \). El menor autovalor es \( 2 \), por lo que \( \nabla^2 f(x) \succeq 2I \) para todo \( (x_1, x_2) \). Por lo tanto, \( f \) es fuertemente convexa con parámetro \( \mu = 2 \).
%   \end{enumerate}
\end{ex}

\section{Convergencia débil}

\textbf{Definición:}
Una sucesión \( \{v_n\} \) en un espacio de Hilbert \( V \) \textbf{converge débilmente} a \( v \in V \) (se denota \( v_n \rightharpoonup v \)) si para todo \( w \in V \) se cumple:
\[
\langle v_n, w \rangle_V \to \langle v, w \rangle \quad \text{cuando} \quad n \to \infty.
\]

\vspace{1em}

% \begin{ex}
%   En el espacio \( \ell^2 \) (secuencias cuadrado-summables), considera la sucesión \( \{e_n\} \) donde \( e_n = (0, \dots, 0, 1, 0, \dots) \) (el 1 está en la posición \( n \)-ésima).
%   \begin{enumerate}
%     \item Demuestra que \( e_n \rightharpoonup 0 \).
%     \item ¿Converge \( \{e_n\} \) fuertemente a 0? Justifica tu respuesta.
%   \end{enumerate}

%   \textbf{Corrección:}
%   \begin{enumerate}
%     \item Sea \( y = (y_1, y_2, \dots) \in \ell^2 \). Entonces \( \langle e_n, y \rangle = y_n \). Como \( y \in \ell^2 \), se tiene \( y_n \to 0 \) cuando \( n \to \infty \). Por lo tanto, \( \langle e_n, y \rangle \to 0 = \langle 0, y \rangle \), lo que implica \( e_n \rightharpoonup 0 \).
%     \item No, \( \{e_n\} \) no converge fuertemente a 0 porque \( \|e_n - 0\|_2 = \|e_n\|_2 = 1 \) para todo \( n \), y no tiende a 0.
%   \end{enumerate}
% \end{ex}

\begin{ex}
  En \( \ell^2 \), sea \( x_n = \left(1, \frac{1}{2}, \frac{1}{3}, \dots, \frac{1}{n}, 0, 0, \dots \right) \).
  \begin{enumerate}
    \item Demuestra que \( x_n \to x \) fuertemente en \( \ell^2 \), donde \( x = \left(1, \frac{1}{2}, \frac{1}{3}, \dots \right) \).
    \item ¿Converge \( \{x_n\} \) débilmente a \( x \)? Justifica tu respuesta.
  \end{enumerate}

%   \textbf{Corrección:}
%   \begin{enumerate}
%     \item Se tiene \( \|x_n - x\|_2^2 = \sum_{k=n+1}^\infty \frac{1}{k^2} \to 0 \) cuando \( n \to \infty \), ya que la serie \( \sum_{k=1}^\infty \frac{1}{k^2} \) converge. Por lo tanto, \( x_n \to x \) fuertemente.
%     \item Sí, porque la convergencia fuerte implica convergencia débil.
%   \end{enumerate}
\end{ex}

% \begin{ex}
%   En \( \ell^2 \), considera la sucesión \( x_n = \left(\frac{1}{n}, \frac{1}{n}, \frac{1}{n}, \dots \right) \).
%   \begin{enumerate}
%     \item Demuestra que \( x_n \to 0 \) fuertemente.
%     \item ¿Converge \( \{x_n\} \) débilmente a 0?
%   \end{enumerate}

%   \textbf{Corrección:}
%   \begin{enumerate}
%     \item Se tiene \( \|x_n\|_2^2 = \sum_{k=1}^\infty \left(\frac{1}{n}\right)^2 \), pero esta serie diverge (es infinita). Por lo tanto, \( x_n \notin \ell^2 \) para ningún \( n \). \textbf{Nota:} Este ejercicio muestra que la sucesión no está bien definida en \( \ell^2 \). \textit{Corrección del ejercicio:} Considera \( x_n = \left(\frac{1}{n}, \frac{1}{n^2}, \frac{1}{n^3}, \dots \right) \). Entonces \( \|x_n\|_2^2 = \sum_{k=1}^\infty \frac{1}{n^{2k}} = \frac{1/n^2}{1 - 1/n^2} \to 0 \), por lo que \( x_n \to 0 \) fuertemente, y por lo tanto, débilmente.
%   \end{enumerate}
% \end{ex}

% \begin{ex}
%   En \( \ell^2 \), sea \( x_n = \left(1, 0, 0, \dots \right) \) para todo \( n \) (es decir, \( x_n = e_1 \) para todo \( n \)).
%   \begin{enumerate}
%     \item ¿Converge \( \{x_n\} \) débilmente? ¿A qué?
%     \item ¿Converge \( \{x_n\} \) fuertemente?
%   \end{enumerate}

%   \textbf{Corrección:}
%   \begin{enumerate}
%     \item Sí, \( x_n \rightharpoonup e_1 \), ya que \( \langle x_n, y \rangle = y_1 = \langle e_1, y \rangle \) para todo \( y \in \ell^2 \).
%     \item Sí, \( x_n \to e_1 \) fuertemente porque \( \|x_n - e_1\|_2 = 0 \) para todo \( n \).
%   \end{enumerate}
% \end{ex}

\begin{ex}
  En \( L^2(0,1) \), considera la sucesión de funciones \( f_n(x) = x^n \) para \( x \in [0,1] \).
  \begin{enumerate}
    \item Demuestra que \( f_n \rightharpoonup 0 \) en \( L^2(0,1) \).
    \item ¿Converge \( \{f_n\} \) fuertemente a 0?
  \end{enumerate}

%   \textbf{Corrección:}
%   \begin{enumerate}
%     \item Sea \( g \in L^2(0,1) \). Entonces \( \langle f_n, g \rangle = \int_0^1 x^n g(x) \, dx \). Como \( |x^n g(x)| \leq |g(x)| \) y \( x^n g(x) \to 0 \) puntualmente en \( [0,1) \), por el teorema de convergencia dominada de Lebesgue, \( \langle f_n, g \rangle \to 0 \). Por lo tanto, \( f_n \rightharpoonup 0 \).
%     \item No, porque \( \|f_n\|_2^2 = \int_0^1 x^{2n} \, dx = \frac{1}{2n+1} \to 0 \), pero la convergencia es fuerte en este caso. \textbf{Nota:} En realidad, \( \|f_n\|_2 \to 0 \), por lo que \( f_n \to 0 \) fuertemente. \textit{Corrección del ejercicio:} Considera \( f_n(x) = \sin(n \pi x) \). Entonces \( \|f_n\|_2^2 = \int_0^1 \sin^2(n \pi x) \, dx = \frac{1}{2} \), por lo que \( \{f_n\} \) no converge fuertemente a 0, pero \( \langle f_n, g \rangle = \int_0^1 \sin(n \pi x) g(x) \, dx \to 0 \) para todo \( g \in L^2(0,1) \) (por el lemma de Riemann-Lebesgue). Por lo tanto, \( f_n \rightharpoonup 0 \).
%   \end{enumerate}
\end{ex}

\begin{ex}
  En \( L^2(0,1) \), considera la sucesión \( f_n(x) = \sqrt{n} \chi_{[0,1/n]} \).
  \begin{enumerate}
    \item Demuestra que \( \|f_n\|_2 = 1 \) para todo \( n \).
    \item Demuestra que \( f_n \rightharpoonup 0 \).
    \item ¿Converge \( \{f_n\} \) fuertemente a 0?
  \end{enumerate}

%   \textbf{Corrección:}
%   \begin{enumerate}
%     \item Se tiene \( \|f_n\|_2^2 = \int_0^{1/n} (\sqrt{n})^2 \, dx = \int_0^{1/n} n \, dx = n \cdot \frac{1}{n} = 1 \), por lo que \( \|f_n\|_2 = 1 \).
%     \item Sea \( g \in L^2(0,1) \). Entonces, por la desigualdad de Cauchy-Schwarz,
%       \[
%       |\langle f_n, g \rangle| = \left| \sqrt{n} \int_0^{1/n} g(x) \, dx \right| \leq \sqrt{n} \left( \int_0^{1/n} 1^2 \, dx \right)^{1/2} \left( \int_0^{1/n} |g(x)|^2 \, dx \right)^{1/2} = \|g\|_{L^2[0,1/n]}.
%       \]
%       Como \( g \in L^2(0,1) \), se cumple \( \int_0^{1/n} |g(x)|^2 \, dx \to 0 \) cuando \( n \to \infty \) (pues la medida de \( [0,1/n] \) tiende a 0). Por lo tanto, \( \|g\|_{L^2[0,1/n]} \to 0 \), y así \( \langle f_n, g \rangle \to 0 \). Esto implica que \( f_n \rightharpoonup 0 \).
%     \item No, porque \( \|f_n - 0\|_2 = 1 \) para todo \( n \), por lo que no tiende a 0.
%   \end{enumerate}
\end{ex}

\begin{ex}
  En \( L^2(0,1) \), considera la sucesión \( f_n(x) = \cos(2 \pi n x) \).
  \begin{enumerate}
    \item Demuestra que \( f_n \rightharpoonup 0 \) en \( L^2(0,1) \).
    \item ¿Converge \( \{f_n\} \) fuertemente a 0?
  \end{enumerate}

%   \textbf{Corrección:}
%   \begin{enumerate}
%     \item Sea \( g \in L^2(0,1) \). Entonces \( \langle f_n, g \rangle = \int_0^1 \cos(2 \pi n x) g(x) \, dx \). Por el lemma de Riemann-Lebesgue, esta integral tiende a 0 cuando \( n \to \infty \). Por lo tanto, \( f_n \rightharpoonup 0 \).
%     \item No, porque \( \|f_n\|_2^2 = \int_0^1 \cos^2(2 \pi n x) \, dx = \frac{1}{2} \) para todo \( n \), por lo que \( \|f_n\|_2 = \frac{\sqrt{2}}{2} \) no tiende a 0.
%   \end{enumerate}
\end{ex}

\section{Mínimos}

\begin{ex}
Sea $V=\ell_2$ y
\[
\mathcal{J}(v):=
\left ( \| v \|^2_V - 1 
\right )^2 +
\sum_{i\geq1} \frac{x_i^2}{i}.
\]
Mostrar que $\mathcal{J}$ no tiene ningún punto de mínimo en $V$.
\end{ex}

\begin{ex} %[M\'inimo alcanzado, caso f\'acil]
Sea $V=\ell_2$ y
\[
\mathcal{J}(v):=\sum_{i\geq 1}\frac{x_i^2}{i},\qquad v=(x_1,x_2,\dots).
\]
Mostrar que $\inf_V\mathcal J=0$, que es alcanzado, y describir el conjunto de
minimizantes.
\end{ex}

% \begin{sol}
% Como cada t\'ermino es $\ge 0$, se tiene $\mathcal J\ge 0$, y $\mathcal J(0)=0$,
% as\'{\i} que $\inf\mathcal J=0$ y $0$ es minimizante. Si $\mathcal J(v)=0$,
% cada t\'ermino $x_i^2/i$ es nulo, luego $x_i=0$ para todo $i$ y $v=0$. El
% m\'{\i}nimo $0$ se alcanza \emph{\'unicamente} en $v=0$.
% \end{sol}

\begin{ex}[\'Infimo no alcanzado en la esfera]
Sea $V=\ell_2$ y
\[
\mathcal{J}(v):=\sum_{i\geq 1}\frac{x_i^2}{i}
\qquad\text{restringida a }\|v\|_{\ell_2}=1.
\]
Mostrar que el \'infimo de $\mathcal J$ sobre la esfera unidad es $0$, pero que
ning\'un $v$ con $\|v\|_{\ell_2}=1$ lo alcanza.
\end{ex}

% \begin{sol}
% Para $e_n$ el $n$-\'esimo vector de la base can\'onica,
% $\mathcal J(e_n)=\frac1n\to 0$, luego $\inf\le 0$; y como $\mathcal J>0$ salvo en
% $0$ (que no est\'a en la esfera), $\inf=0$.
% Si $v$ estuviera en la esfera con $\mathcal J(v)=0$, el ejercicio anterior da
% $v=0$, contradicci\'on con $\norm{v}=1$. El \'infimo no se alcanza: la esfera de
% $\ell_2$ no es compacta.
% \end{sol}

\begin{center}
    $\star$
\end{center}


\chapter{Newton semi-suave}

\section{L\'imites inferiores}

\begin{ex}
  Calcula el límite inferior de las siguientes sucesiones:
  \begin{enumerate}
    \item \( a_n = 1 + \frac{(-1)^n}{n} \),
    \item \( b_n = \frac{n + (-1)^n (n+1)}{n} \).
  \end{enumerate}

%   \textbf{Corrección:}
%   \begin{enumerate}
%     \item Para \( a_n = 1 + \frac{(-1)^n}{n} \):
%       \begin{itemize}
%         \item Si \( n \) es par, \( a_n = 1 + \frac{1}{n} \).
%         \item Si \( n \) es impar, \( a_n = 1 - \frac{1}{n} \).
%       \end{itemize}
%       Los términos \( a_n \) para \( n \) par son \( 1 + \frac{1}{2}, 1 + \frac{1}{4}, 1 + \frac{1}{6}, \dots \), que convergen a \( 1 \) por arriba.
%       Los términos \( a_n \) para \( n \) impar son \( 1 - \frac{1}{1}, 1 - \frac{1}{3}, 1 - \frac{1}{5}, \dots \), que convergen a \( 1 \) por abajo.
%       Por lo tanto, el conjunto de puntos límite de \( \{a_n\} \) es \( \{1\} \), y así:
%       \[
%       \liminf_{n \to \infty} a_n = 1.
%       \]

%     \item Para \( b_n = \frac{n + (-1)^n (n+1)}{n} \):
%       \begin{itemize}
%         \item Si \( n \) es par, \( b_n = \frac{n + (n+1)}{n} = \frac{2n + 1}{n} = 2 + \frac{1}{n} \).
%         \item Si \( n \) es impar, \( b_n = \frac{n - (n+1)}{n} = \frac{-1}{n} \).
%       \end{itemize}
%       Los términos \( b_n \) para \( n \) par convergen a \( 2 \), y los términos para \( n \) impar convergen a \( 0 \).
%       El conjunto de puntos límite de \( \{b_n\} \) es \( \{0, 2\} \), por lo que:
%       \[
%       \liminf_{n \to \infty} b_n = \min\{0, 2\} = 0.
%       \]
%   \end{enumerate}
\end{ex}

\begin{ex}
  Sea \( \{a_n\} \) una sucesión acotada. Demuestra que:
  \[
  \liminf_{n \to \infty} a_n \leq \limsup_{n \to \infty} a_n.
  \]
  Además, si \( \liminf_{n \to \infty} a_n = \limsup_{n \to \infty} a_n = L \), demuestra que \( \{a_n\} \) converge a \( L \).

%   \textbf{Corrección:}
%   \begin{itemize}
%     \item Por definición, \( \liminf_{n \to \infty} a_n = \lim_{n \to \infty} \left( \inf_{k \geq n} a_k \right) \) y \( \limsup_{n \to \infty} a_n = \lim_{n \to \infty} \left( \sup_{k \geq n} a_k \right) \).
%       Para todo \( n \), se cumple \( \inf_{k \geq n} a_k \leq \sup_{k \geq n} a_k \), ya que el ínfimo de un conjunto es menor o igual que su supremo.
%       Como el límite preserva desigualdades, tomamos el límite cuando \( n \to \infty \) para obtener:
%       \[
%       \liminf_{n \to \infty} a_n \leq \limsup_{n \to \infty} a_n.
%       \]

%     \item Si \( \liminf_{n \to \infty} a_n = \limsup_{n \to \infty} a_n = L \), entonces para todo \( \varepsilon > 0 \), existen \( N_1, N_2 \in \mathbb{N} \) tales que:
%       \[
%       \inf_{k \geq N_1} a_k > L - \varepsilon \quad \text{y} \quad \sup_{k \geq N_2} a_k < L + \varepsilon.
%       \]
%       Tomando \( N = \max(N_1, N_2) \), para todo \( k \geq N \), se cumple:
%       \[
%       L - \varepsilon < \inf_{k \geq N} a_k \leq a_k \leq \sup_{k \geq N} a_k < L + \varepsilon.
%       \]
%       Por lo tanto, \( |a_k - L| < \varepsilon \) para todo \( k \geq N \), lo que demuestra que \( \{a_n\} \) converge a \( L \).
%   \end{itemize}
\end{ex}

\begin{ex}
  Sea \( \{a_n\} \) una sucesión tal que \( a_n \geq 0 \) para todo \( n \). Demuestra que:
  \[
  \liminf_{n \to \infty} a_n \geq 0.
  \]
  ¿Se cumple el resultado si \( a_n \geq c \) para alguna constante \( c \in \mathbb{R} \)? Justifica tu respuesta.

%   \textbf{Corrección:}
%   \begin{itemize}
%     \item Como \( a_n \geq 0 \) para todo \( n \), el ínfimo de cualquier subconjunto de \( \{a_n\} \) también es \( \geq 0 \). Por lo tanto, \( \inf_{k \geq n} a_k \geq 0 \) para todo \( n \), y así:
%       \[
%       \liminf_{n \to \infty} a_n = \lim_{n \to \infty} \left( \inf_{k \geq n} a_k \right) \geq 0.
%       \]

%     \item Sí, el resultado se generaliza. Si \( a_n \geq c \) para todo \( n \), entonces \( \inf_{k \geq n} a_k \geq c \) para todo \( n \), y por lo tanto:
%       \[
%       \liminf_{n \to \infty} a_n \geq c.
%       \]
%   \end{itemize}
\end{ex}

\begin{ex}
  Calcula el límite inferior de las siguientes sucesiones:
  \begin{enumerate}
    \item \( c_n = \begin{cases}
      n & \text{si } n \text{ es impar}, \\
      \frac{1}{n} & \text{si } n \text{ es par}.
    \end{cases} \)
    \item \( d_n = \sin\left(\frac{n \pi}{2}\right) \).
  \end{enumerate}

%   \textbf{Corrección:}
%   \begin{enumerate}
%     \item Para \( c_n \):
%       \begin{itemize}
%         \item Si \( n \) es impar, \( c_n = n \to \infty \).
%         \item Si \( n \) es par, \( c_n = \frac{1}{n} \to 0 \).
%       \end{itemize}
%       El conjunto de puntos límite de \( \{c_n\} \) es \( \{0\} \cup \{\infty\} \), pero como \( \infty \) no es un punto límite finito, el único punto límite finito es \( 0 \).
%       Por lo tanto:
%       \[
%       \liminf_{n \to \infty} c_n = 0.
%       \]

%     \item Para \( d_n = \sin\left(\frac{n \pi}{2}\right) \):
%       \begin{itemize}
%         \item Si \( n = 4k \), \( d_n = \sin(2k\pi) = 0 \).
%         \item Si \( n = 4k+1 \), \( d_n = \sin\left(\frac{(4k+1)\pi}{2}\right) = \sin\left(2k\pi + \frac{\pi}{2}\right) = 1 \).
%         \item Si \( n = 4k+2 \), \( d_n = \sin\left(\frac{(4k+2)\pi}{2}\right) = \sin\left(2k\pi + \pi\right) = 0 \).
%         \item Si \( n = 4k+3 \), \( d_n = \sin\left(\frac{(4k+3)\pi}{2}\right) = \sin\left(2k\pi + \frac{3\pi}{2}\right) = -1 \).
%       \end{itemize}
%       Los valores de la sucesión son \( 0, 1, 0, -1, 0, 1, 0, -1, \dots \), por lo que el conjunto de puntos límite es \( \{-1, 0, 1\} \).
%       Por lo tanto:
%       \[
%       \liminf_{n \to \infty} d_n = \min\{-1, 0, 1\} = -1.
%       \]
%   \end{enumerate}
\end{ex}

\begin{ex}
  Sea \( \{a_n\} \) una sucesión tal que \( \liminf_{n \to \infty} a_n = L \). Demuestra que existe una subsucesión \( \{a_{n_k}\} \) tal que \( a_{n_k} \to L \).

%   \textbf{Corrección:}
%   Por definición, \( \liminf_{n \to \infty} a_n = \lim_{n \to \infty} \left( \inf_{k \geq n} a_k \right) = L \).
%   Para cada \( n \), sea \( m_n = \inf_{k \geq n} a_k \). Entonces \( m_n \to L \) cuando \( n \to \infty \).

%   Para cada \( n \), existe \( k_n \geq n \) tal que:
%   \[
%   m_n \leq a_{k_n} \leq m_n + \frac{1}{n}.
%   \]
%   Esto es posible porque \( m_n \) es el ínfimo de \( \{a_k \mid k \geq n\} \), por lo que para todo \( \varepsilon > 0 \) (en este caso, \( \varepsilon = 1/n \)), existe \( k_n \geq n \) tal que \( a_{k_n} \leq m_n + \varepsilon \).

%   Consideremos la subsucesión \( \{a_{k_n}\} \). Entonces:
%   \[
%   |a_{k_n} - L| \leq |a_{k_n} - m_n| + |m_n - L| \leq \frac{1}{n} + |m_n - L|.
%   \]
%   Como \( m_n \to L \) y \( \frac{1}{n} \to 0 \), se sigue que \( a_{k_n} \to L \).
%   Además, \( \{k_n\} \) es estrictamente creciente (ya que \( k_n \geq n \) y \( n \) crece), por lo que \( \{a_{k_n}\} \) es efectivamente una subsucesión de \( \{a_n\} \).
\end{ex}

\section{Semi-continuidad inferior}

\begin{ex}
  Determina si las siguientes funciones \( f: \mathbb{R} \to \mathbb{R} \) son semicontinuas inferiormente (s.c.i.):
  \begin{enumerate}
    \item \( f(x) = |x| \).
    \item \( f(x) = \begin{cases}
      -1 & \text{si } x < 0, \\
      1 & \text{si } x \geq 0.
    \end{cases} \)
  \end{enumerate}

%   \textbf{Corrección:}
%   \begin{enumerate}
%     \item \textbf{\( f(x) = |x| \) es s.c.i.}
%       La función \( f(x) = |x| \) es continua en \( \mathbb{R} \). Como toda función continua es s.c.i., \( f \) es s.c.i.

%     \item \textbf{La función no es s.c.i. en \( x = 0 \).}
%       Para verificar la semicontinuidad inferior en \( x = 0 \), consideramos una sucesión \( \{x_n\} \) tal que \( x_n \to 0 \) y \( x_n < 0 \) para todo \( n \). Entonces \( f(x_n) = -1 \) para todo \( n \), y por lo tanto:
%       \[
%       \liminf_{n \to \infty} f(x_n) = -1 < f(0) = 1.
%       \]
%       Como existe una sucesión convergente a \( 0 \) para la cual \( \liminf_{n \to \infty} f(x_n) < f(0) \), la función no es s.c.i. en \( x = 0 \).
%   \end{enumerate}
\end{ex}

\begin{ex}
  Determina si las siguientes funciones \( f: \mathbb{R} \to \mathbb{R} \) son semicontinuas inferiormente (s.c.i.):
  \begin{enumerate}
    \item \( f(x) = x^2 \).
    \item \( f(x) = \begin{cases}
      x & \text{si } x \leq 0, \\
      x + 1 & \text{si } x > 0.
    \end{cases} \)
  \end{enumerate}

%   \textbf{Corrección:}
%   \begin{enumerate}
%     \item \textbf{\( f(x) = x^2 \) es s.c.i.}
%       La función \( f(x) = x^2 \) es continua en \( \mathbb{R} \), por lo que es s.c.i.

%     \item \textbf{La función no es s.c.i. en \( x = 0 \).}
%       Consideremos la sucesión \( x_n = \frac{1}{n} \), que converge a \( 0 \) por la derecha. Entonces \( f(x_n) = x_n + 1 = \frac{1}{n} + 1 \), y así:
%       \[
%       \liminf_{n \to \infty} f(x_n) = 1 > f(0) = 0.
%       \]
%       Sin embargo, la definición de s.c.i. requiere que \( \liminf_{n \to \infty} f(x_n) \geq f(0) \) para \textbf{toda} sucesión \( \{x_n\} \to 0 \). Aquí, \( \liminf_{n \to \infty} f(x_n) = 1 \geq 0 = f(0) \), por lo que la condición se cumple para esta sucesión.
%       \textit{Corrección:} Consideremos la sucesión \( x_n = -\frac{1}{n} \), que converge a \( 0 \) por la izquierda. Entonces \( f(x_n) = x_n = -\frac{1}{n} \), y así:
%       \[
%       \liminf_{n \to \infty} f(x_n) = 0 = f(0).
%       \]
%       Para cualquier sucesión \( \{x_n\} \to 0 \):
%       \begin{itemize}
%         \item Si \( x_n \leq 0 \) para todo \( n \), entonces \( f(x_n) = x_n \to 0 = f(0) \), por lo que \( \liminf_{n \to \infty} f(x_n) = 0 \geq f(0) \).
%         \item Si \( x_n > 0 \) para infinitos \( n \), entonces \( f(x_n) = x_n + 1 \to 1 \), por lo que \( \liminf_{n \to \infty} f(x_n) \geq 0 = f(0) \).
%       \end{itemize}
%       Por lo tanto, \( \liminf_{n \to \infty} f(x_n) \geq f(0) \) para toda sucesión \( \{x_n\} \to 0 \), y la función \textbf{es s.c.i. en \( x = 0 \)}.
%       \textit{Nota:} La función es s.c.i. en todo \( \mathbb{R} \), ya que es continua en \( \mathbb{R} \setminus \{0\} \) y en \( x = 0 \) cumple la condición de s.c.i.
%   \end{enumerate}
\end{ex}

% \begin{ex}
%   Demuestra que la función \( f: \mathbb{R} \to \mathbb{R} \) definida por
%   \[
%   f(x) = \begin{cases}
%     0 & \text{si } x \leq 0, \\
%     x & \text{si } x > 0,
%   \end{cases}
%   \]
%   es semicontinua inferiormente en \( \mathbb{R} \).

%   \textbf{Corrección:}
%   La función \( f \) es continua en \( \mathbb{R} \setminus \{0\} \), por lo que solo necesitamos verificar la semicontinuidad inferior en \( x = 0 \).
%   Sea \( \{x_n\} \) una sucesión tal que \( x_n \to 0 \). Consideramos dos casos:
%   \begin{itemize}
%     \item Si \( x_n \leq 0 \) para todo \( n \), entonces \( f(x_n) = 0 \) para todo \( n \), y así:
%       \[
%       \liminf_{n \to \infty} f(x_n) = 0 = f(0).
%       \]
%     \item Si \( x_n > 0 \) para infinitos \( n \), entonces para esos \( n \), \( f(x_n) = x_n \to 0 \), por lo que:
%       \[
%       \liminf_{n \to \infty} f(x_n) \geq 0 = f(0).
%       \]
%   \end{itemize}
%   En ambos casos, \( \liminf_{n \to \infty} f(x_n) \geq f(0) \), lo que demuestra que \( f \) es s.c.i. en \( x = 0 \).
% \end{ex}

\begin{ex}
  Sea \( f: \mathbb{R} \to \mathbb{R} \) una función s.c.i. y sea \( g: \mathbb{R} \to \mathbb{R} \) definida por \( g(x) = f(x) + c \), donde \( c \in \mathbb{R} \) es una constante. Demuestra que \( g \) es s.c.i.

%   \textbf{Corrección:}
%   Sea \( \{x_n\} \) una sucesión tal que \( x_n \to x \). Como \( f \) es s.c.i., se cumple:
%   \[
%   \liminf_{n \to \infty} f(x_n) \geq f(x).
%   \]
%   Entonces:
%   \[
%   \liminf_{n \to \infty} g(x_n) = \liminf_{n \to \infty} (f(x_n) + c) = \liminf_{n \to \infty} f(x_n) + c \geq f(x) + c = g(x).
%   \]
%   Por lo tanto, \( g \) es s.c.i.
\end{ex}

\begin{ex}
  Sea \( f, g: \mathbb{R} \to \mathbb{R} \) funciones s.c.i. Demuestra que \( h(x) = \max(f(x), g(x)) \) es s.c.i.

%   \textbf{Corrección:}
%   Sea \( \{x_n\} \) una sucesión tal que \( x_n \to x \). Como \( f \) y \( g \) son s.c.i., se cumple:
%   \[
%   \liminf_{n \to \infty} f(x_n) \geq f(x) \quad \text{y} \quad \liminf_{n \to \infty} g(x_n) \geq g(x).
%   \]
%   Para cada \( n \), se tiene \( h(x_n) = \max(f(x_n), g(x_n)) \geq f(x_n) \) y \( h(x_n) \geq g(x_n) \). Por lo tanto:
%   \[
%   \liminf_{n \to \infty} h(x_n) \geq \liminf_{n \to \infty} f(x_n) \geq f(x) \quad \text{y} \quad \liminf_{n \to \infty} h(x_n) \geq \liminf_{n \to \infty} g(x_n) \geq g(x).
%   \]
%   Esto implica:
%   \[
%   \liminf_{n \to \infty} h(x_n) \geq \max(f(x), g(x)) = h(x).
%   \]
%   Por lo tanto, \( h \) es s.c.i.
\end{ex}

% \chapter{Desigualdades variacionales}

% \dots

\section{Subdiferencial}

% ============================================================
% Ejercicios de análisis convexo: la subdiferencial
% 10 ejercicios con soluciones detalladas, en español.
% Requiere: amsmath, amssymb, mathtools, amsthm.
% Entornos: \begin{ex} ... \end{ex}  (como en el modelo).
% ============================================================

% ---------- Ejercicio 1 ----------
\begin{ex}
Sea \( f(x) = |x| \). Calcula \( \partial f(x) \) para todo \( x \in \mathbb{R} \), y deduce el conjunto de los minimizadores de \( f \).
\end{ex}

% \begin{sol}
% Por definición, \( s \in \partial f(x) \) si y solo si
% \[
% f(y) \geq f(x) + s (y - x) \quad \text{para todo } y \in \mathbb{R}.
% \]

% \emph{Caso 1: \( x > 0 \).} Entonces \( f \) es diferenciable en \( x \) con \( f'(x) = 1 \). Como toda función convexa diferenciable satisface \( \partial f(x) = \{f'(x)\} \), tenemos \( \partial f(x) = \{1\} \).

% Verifiquemoslo directamente: para todo \( y \),
% \[
% |y| \geq |x| + (y - x)
% \]
% pues si \( y \geq 0 \) la desigualdad es la identidad \( |y| = y \geq x + y - x = |x| + (y-x) \)... En efecto, para \( y \geq 0 \): \( |y| - (y-x) = x = |x| \), hay igualdad; para \( y < 0 \): \( |y| - (y - x) = -2y + x > x = |x| \). Así \( s = 1 \in \partial f(x) \). Si \( s > 1 \), tomando \( y \to +\infty \) se viola la desigualdad; si \( s < 1 \), tomando \( y \to -\infty \) también. Luego \( \partial f(x) = \{1\} \).

% \emph{Caso 2: \( x < 0 \).} Simétricamente, \( \partial f(x) = \{-1\} \).

% \emph{Caso 3: \( x = 0 \).} La condición es \( |y| \geq s\, y \) para todo \( y \). Si \( y > 0 \): \( s \leq 1 \). Si \( y < 0 \): \( s \geq -1 \). Por lo tanto
% \[
% \partial f(0) = [-1, 1].
% \]
% Obsérvese que coincide con \( [f'_-(0), f'_+(0)] = [-1, 1] \), como corresponde en dimensión 1.

% \emph{Minimizadores.} Por el criterio de optimalidad, \( x \) minimiza \( f \) si y solo si \( 0 \in \partial f(x) \). Esto ocurre únicamente en \( x = 0 \), donde \( 0 \in [-1,1] \). El único minimizador es \( x = 0 \), con \( f(0) = 0 \).
% \end{sol}


% ---------- Ejercicio 2 ----------
\begin{ex}
Sea \( f(x) = x^2 \) y \( g(x) = |x| \). Calcula \( \partial (f + g)(x) \). % y verifica el resultado con la regla de la suma.
\end{ex}

% Ambas funciones son convexas y continuas en \( \mathbb{R} \) (condición de calificación satisfecha: todo el espacio), así que la regla de la suma da
% \[
% \partial (f + g)(x) = \partial f(x) + \partial g(x) = \{2x\} + \partial |x|(x).
% \]

% \emph{Fuera del origen:} \( \partial |x|(x) = \{\operatorname{sgn}(x)\} \), luego
% \[
% \partial(f+g)(x) = \{ 2x + \operatorname{sgn}(x) \}, \qquad x \neq 0.
% \]

% \emph{En el origen:} \( \partial(f+g)(0) = \{0\} + [-1,1] = [-1, 1] \).

% \emph{Verificación directa en \( x = 0 \):} \( (f+g)(y) = y^2 + |y| \), y \( s \in \partial(f+g)(0) \) sii
% \[
% y^2 + |y| \geq s\, y \quad \text{para todo } y.
% \]
% Para \( y > 0 \): \( y + 1 \geq s \), o sea \( s \leq 1 \). Para \( y < 0 \): dividiendo por \( y < 0 \), \( s \leq y + 1 \) no—calculamos bien: \( y^2 - y \geq s y \) da, dividiendo por \( y < 0 \), \( s \geq y - 1 \) para todo \( y < 0 \), es decir \( s \geq -1 \). Conclusión: \( s \in [-1,1] \), que coincide con la regla de la suma.

\emph{Observación.} El mínimo de \( f+g \): \( 0 \in \partial(f+g)(x) \) exige \( 2x + \operatorname{sgn}(x) = 0 \) (imposible para \( x \neq 0 \)) o \( x = 0 \) (donde \( 0 \in [-1,1] \)). El mínimo se alcanza en \( x = 0 \). Es el primer contacto con la estructura de los problemas regularizados tipo LASSO.

% ---------- Ejercicio 3 ----------
% \begin{ex}
% Sea \( K = [a, b] \subset \mathbb{R} \) y sea \( \delta_K \) su función indicadora:
% \[
% \delta_K(x) = \begin{cases} 0 & \text{si } x \in K, \\ +\infty & \text{si no.} \end{cases}
% \]
% Calcula \( \partial \delta_K(x) \) para todo \( x \in \mathbb{R} \).
% \end{ex}

% \begin{sol}
% Primero, \( \delta_K \) es convexa y s.c.i. pues \( K \) es convexo y cerrado.

% \emph{Paso 1: interior.} Si \( x \in (a,b) \), \( \delta_K \) vale \( 0 \) en un entorno de \( x \), luego es diferenciable ahí con derivada nula y \( \partial \delta_K(x) = \{0\} \). Verificación: \( 0 \) sirve pues \( \delta_K(y) \geq 0 = \delta_K(x) \) para todo \( y \). Y si \( s \neq 0 \), tomando \( y = x + \varepsilon \operatorname{sgn}(s) \in K \) (posible pues \( x \) es interior) se obtiene \( 0 \geq s \varepsilon |s|/|s| > 0 \), contradicción.

% \emph{Paso 2: borde izquierdo \( x = a \).} \( s \in \partial \delta_K(a) \) sii para todo \( y \in [a,b] \):
% \[
% 0 \geq s (y - a).
% \]
% Como \( y - a \geq 0 \), esto exige \( s \leq 0 \). Luego
% \[
% \partial \delta_K(a) = (-\infty, 0] = N_K(a),
% \]
% el cono normal a \( K \) en \( a \).

% \emph{Paso 3: borde derecho \( x = b \).} Análogamente \( s(y - b) \leq 0 \) para \( y \leq b \) exige \( s \geq 0 \):
% \[
% \partial \delta_K(b) = [0, +\infty) = N_K(b).
% \]

% \emph{Fuera de \( K \):} \( \partial \delta_K(x) = \varnothing \) (la función vale \( +\infty \) ahí; el epígrafo no tiene puntos por encima de \( (x, \delta_K(x)) \)).

% \emph{Moraleja.} \( \partial \delta_K = N_K \): la subdiferencial de la indicadora es el cono normal. Es la puerta de entrada a KKT: minimizar \( f \) sobre \( K \) es minimizar \( f + \delta_K \), y \( 0 \in \partial f(\bar x) + N_K(\bar x) \).
% \end{sol}

% ---------- Ejercicio 4 ----------
\begin{ex}
Calcula \( \partial f \) de la función de ReLU:
\[
f(x) = \operatorname{ReLU}(x) = \max(0, x).
\]
¿En qué puntos existe \( \nabla f \)? ¿Por qué no es una contradicción con \( \partial f(x) \) intervalo no degenerado?
\end{ex}

% Como máximo de dos funciones convexas afines, \( f \) es convexa.

% \emph{Caso \( x < 0 \):} \( f \equiv 0 \) cerca de \( x \), luego \( \partial f(x) = \{0\} \).

% \emph{Caso \( x > 0 \):} \( f(x) = x \) cerca de \( x \), luego \( \partial f(x) = \{1\} \).

% \emph{Caso \( x = 0 \):} \( s \in \partial f(0) \) sii
% \[
% \max(0, y) \geq s\, y \quad \text{para todo } y \in \mathbb{R}.
% \]
% Para \( y > 0 \): \( 1 \geq s \). Para \( y < 0 \): \( 0 \geq s y \), es decir \( s \geq 0 \). Por lo tanto
% \[
% \partial f(0) = [0, 1].
% \]
% (Es de nuevo \( [f'_-(0), f'_+(0)] = [0,1] \).)

% \emph{Diferenciabilidad.} \( f \) es diferenciable (clásicamente) en \( \mathbb{R} \setminus \{0\} \), con \( \nabla f = \mathbb{1}_{x > 0} \). En el origen no existe la derivada clásica, pero la subdiferencial es no vacía (\( [0,1] \)).

% \emph{No hay contradicción.} La implicación válida es
% \[
% f \text{ diferenciable en } x \;\Longrightarrow\; \partial f(x) = \{\nabla f(x)\},
% \]
% pero la recíproca es \emph{falsa}: en dimensión 1, \( \partial f(x) \) puede ser un intervalo no degenerado aunque no exista la derivada. Error típico: creer que \( \partial f(x) \) singleton implica diferenciabilidad.

% ---------- Ejercicio 5 ----------
\begin{ex}
Sea \( A \in \mathbb{R}^{m \times n} \), \( b \in \mathbb{R}^m \), y sea
\[
f(x) = \| A x - b \|_2.
\]
Calcula \( \partial f(x) \) para todo \( x \in \mathbb{R}^n \).
\end{ex}

% \begin{sol}
% Escribimos \( f = g \circ F \) con \( F(x) = Ax - b \) afín y \( g(z) = \|z\|_2 \).

% \emph{Caso 1: \( A x \neq b \).} Entonces \( g \) es diferenciable en \( F(x) \) con
% \[
% \nabla g(z) = \frac{z}{\|z\|_2},
% \]
% y la regla de composición con función afín da
% \[
% \partial f(x) = A^\top \partial g(Ax - b) = \left\{ \frac{A^\top (A x - b)}{\|A x - b\|_2} \right\}.
% \]

% \emph{Caso 2: \( A x = b \).} Entonces \( f(x) = 0 \) es el mínimo global ( \( f \geq 0 \) ). Por la regla de composición,
% \[
% \partial f(x) = A^\top \partial g(0) = A^\top B_2,
% \]
% donde \( B_2 = \{ z \in \mathbb{R}^m : \|z\|_2 \leq 1 \} \) es la bola unitaria cerrada (subdiferencial de la norma en el origen: es su bola dual; en normas euclidianas, la propia bola). En efecto, \( s \in \partial f(x) \) sii
% \[
% \|Ay - b\|_2 \geq \langle s, y - x \rangle \quad \forall y,
% \]
% y escribiendo \( s = A^\top z \) con \( z \in B_2 \), se sigue de \( \|Ay - b\|_2 = \|Ay - b\|_2 - \|Ax-b\|_2 \geq \langle z, Ay - b \rangle = \langle A^\top z, y - x\rangle \) (Cauchy--Schwarz), que vale para todo \( z \in B_2 \).

% \emph{Minimizadores.} \( 0 \in \partial f(x) \) sii \( 0 \in A^\top B_2 \) sii \( Ax = b \) tiene solución (si \( x \) resuelve el sistema, \( 0 = A^\top \cdot 0 \) con \( 0 \in B_2 \)). Conforme: el mínimo de \( f \), si el sistema es compatible, vale \( 0 \) y se alcanza exactamente en las soluciones.
% \end{sol}

% % ---------- Ejercicio 6 ----------
% \begin{ex}
% Sea \( f(x) = |x - 1| + |x + 1| \), \( x \in \mathbb{R} \). Calcula \( \partial f(x) \) en cada punto y determina el conjunto de minimizadores de \( f \). Interpreta geométricamente.
% \end{ex}

% \begin{sol}
% Por la regla de la suma (todas las funciones continuas en \( \mathbb{R} \)):
% \[
% \partial f(x) = \partial |x-1|(x) + \partial |x+1|(x).
% \]

% Calculamos cada término. \( \partial |x - c|(x) = \{\operatorname{sgn}(x-c)\} \) si \( x \neq c \), y \( = [-1,1] \) si \( x = c \). Por lo tanto:

% \begin{itemize}
% \item \( x < -1 \): \( \partial f(x) = \{-1\} + \{-1\} = \{-2\} \).
% \item \( x = -1 \): \( \partial f(-1) = \{-1\} + [-1,1] = [-2, 0] \).
% \item \( -1 < x < 1 \): \( \partial f(x) = \{1\} + \{-1\} = \{0\} \).
% \item \( x = 1 \): \( \partial f(1) = [0,1] + \{1\} = [1, 2] \).
% \item \( x > 1 \): \( \partial f(x) = \{2\} \).
% \end{itemize}

% \emph{Minimizadores.} Buscamos \( x \) tal que \( 0 \in \partial f(x) \):
% \begin{itemize}
% \item en \( (-1, 1) \), \( \partial f(x) = \{0\} \ni 0 \): minimizan;
% \item en \( x = -1 \), \( 0 \in [-2, 0] \): minimiza;
% \item en \( x = 1 \), \( 0 \in [0, 2] \): minimiza;
% \item fuera de \( [-1,1] \), \( \partial f(x) = \{-2\} \) o \( \{2\} \), que no contienen a \( 0 \).
% \end{itemize}
% Conclusión: el conjunto de minimizadores es exactamente \( [-1, 1] \), y sobre ese segmento \( f(x) = 2 \) constante.

% \emph{Interpretación.} Entre \( -1 \) y \( 1 \) la pendiente izquierda (\( -1 \)) compensa exactamente la pendiente derecha (\( +1 \)): la suma de distancias a los puntos \( \pm 1 \) es constante \( = 2 \). La subdiferencial lo detecta: \( \partial f \supset \{0\} \) en todo el segmento. Es exactamente lo que ocurre, en dimensión superior, con \( f(x) = \|x\|_1 \): el minimizador sobre una recta que corta el interior del "rombo" se aplana sobre una cara entera.
% \end{sol}

% % ---------- Ejercicio 7 ----------
% \begin{ex}
% Sea \( g: \mathbb{R} \to \mathbb{R} \) convexa diferenciable y sea \( f(x) = g(\max(0, x)) \).
% \emph{(a)} Muestra que \( f \) es convexa. \emph{(b)} Calcula \( \partial f(x) \) suponiendo además \( g \) creciente.
% \end{ex}

% \begin{sol}
% \emph{(a)} La función \( x \mapsto \max(0,x) = \operatorname{ReLU}(x) \) es convexa, y \( g \) es convexa y creciente. La composición de una función convexa con una función convexa y creciente es convexa: si \( t \in [0,1] \),
% \[
% f(tx + (1-t)y) = g(\operatorname{ReLU}(tx + (1-t)y)) \leq g( t \operatorname{ReLU}(x) + (1-t)\operatorname{ReLU}(y)) \leq t f(x) + (1-t) f(y),
% \]
% donde usamos la convexidad de ReLU, la monotonía de \( g \) y luego la convexidad de \( g \).

% \emph{(b)} Escribimos \( f = g \circ h \) con \( h = \operatorname{ReLU} \). Por la regla de composición:
% \[
% \partial f(x) = \partial g(h(x)) \cdot \partial h(x) = \{ g'(h(x)) \, s : s \in \partial h(x) \}.
% \]
% (con la convención de que si \( g \) es diferenciable, \( \partial g(z) = \{g'(z)\} \)). Calculamos:

% \begin{itemize}
% \item \( x > 0 \): \( h(x) = x \), \( \partial h(x) = \{1\} \), luego \( \partial f(x) = \{ g'(x) \} \).
% \item \( x < 0 \): \( h(x) = 0 \), \( \partial h(x) = \{0\} \), luego \( \partial f(x) = \{ g'(0) \cdot 0 \} = \{0\} \). Conforme: \( f \equiv g(0) \) constante en \( (-\infty, 0] \).
% \item \( x = 0 \): \( \partial h(0) = [0,1] \), luego
% \[
% \partial f(0) = [0, g'(0)].
% \]
% \end{itemize}

% \emph{Ejemplo (softplus):} con \( g(u) = \log(1 + e^u) \) (creciente, convexa), \( f \) es el softplus suavizado de ReLU y \( \partial f(0) = [0, \tfrac12] \).
% \end{sol}

% % ---------- Ejercicio 8 ----------
% \begin{ex}
% Sea \( K \subset \mathbb{R}^n \) convexo cerrado no vacío, y sea \( f(x) = \operatorname{dist}(x, K) = \min_{y \in K} \|x - y\|_2 \).
% Muestra que \( f \) es convexa y 1-Lipschitz, y calcula \( \partial f(x) \) para \( x \notin K \).
% \end{ex}

% \begin{sol}
% \emph{Convexidad.} \( f \) es ínfimo sobre \( y \in K \) de las funciones \( x \mapsto \|x - y\|_2 \), todas convexas. El ínfimo de funciones convexas en general \emph{no} es convexo, pero aquí el ínfimo se alcanza ( \( K \) cerrado) y podemos usar el criterio del epígrafo o directamente: sean \( x_1, x_2 \in \mathbb{R}^n \), \( t \in [0,1] \), y sean \( p_1, p_2 \in K \) las proyecciones correspondientes. Entonces \( p = t p_1 + (1-t) p_2 \in K \) (convexidad de \( K \)), y
% \[
% f(tx_1 + (1-t)x_2) \leq \| tx_1 + (1-t)x_2 - p \| \leq t \|x_1 - p_1\| + (1-t) \|x_2 - p_2\| = t f(x_1) + (1-t) f(x_2).
% \]

% \emph{1-Lipschitz.} Para \( x_1, x_2 \) y proyección \( p_2 \) de \( x_2 \):
% \[
% f(x_1) \leq \|x_1 - p_2\| \leq \|x_1 - x_2\| + \|x_2 - p_2\| = \|x_1 - x_2\| + f(x_2),
% \]
% y simétricamente; luego \( |f(x_1) - f(x_2)| \leq \|x_1 - x_2\| \).

% \emph{Subdiferencial fuera de \( K \).} Sea \( x \notin K \) y \( p = P_K(x) \) su proyección única. Escribimos
% \[
% f(x) = \|x - p\|_2 = g(F(x)), \qquad F(x) = x - p \ \text{(afín)}, \quad g = \|\cdot\|_2.
% \]
% Como \( x \neq p \), \( g \) es diferenciable en \( F(x) \neq 0 \), y la regla de composición afín da
% \[
% \partial f(x) = \left\{ \frac{x - P_K(x)}{\operatorname{dist}(x, K)} \right\}.
% \]
% Es un vector unitario: el gradiente de la distancia apunta del punto hacia su proyección, ortogonalmente a \( K \) en \( P_K(x) \) (conormalidad de la proyección).

% \emph{Complemento.} En \( x \in K \), \( f(x) = 0 \) es mínimo global y \( 0 \in \partial f(x) \); pero \( \partial f(x) \) puede ser más grande: \( \partial f(x) = N_K(x) \cap B_2 \) cuando \( K \) es convexo cerrado (caso particular: \( K = \{0\} \) da \( \partial \|\cdot\|_2(0) = B_2 \)).
% \end{sol}

% % ---------- Ejercicio 9 ----------
% \begin{ex}
% Sea
% \[
% f(x) = \begin{cases} x \log x & \text{si } x > 0, \\ 0 & \text{si } x = 0, \\ +\infty & \text{si } x < 0. \end{cases}
% \]
% \emph{(a)} Muestra que \( f \) es convexa, s.c.i. y que \( \operatorname{int}(\operatorname{dom} f) = (0, +\infty) \).
% \emph{(b)} Calcula \( \partial f(0) \) y verifica la no-vacuidad en la frontera del dominio.
% \end{ex}

% \begin{sol}
% \emph{(a)} En \( (0, +\infty) \), \( f''(x) = 1/x > 0 \): \( f \) es estrictamente convexa en el interior de su dominio. El dominio \( [0, +\infty) \) es convexo y \( f \) es continua a la derecha en \( 0 \) con \( f(0) = 0 = \lim_{x \to 0^+} x \log x \), así que \( f \) es propia, convexa y s.c.i. (su epígrafo es cerrado). Claramente \( \operatorname{int}(\operatorname{dom} f) = (0, +\infty) \).

% \emph{(b)} \( s \in \partial f(0) \) sii para todo \( y > 0 \):
% \[
% y \log y \geq s\, y, \qquad \text{es decir} \qquad s \leq \log y \quad \forall y > 0.
% \]
% Como \( \log y \to -\infty \) cuando \( y \to 0^+ \), ningún \( s \in \mathbb{R} \) satisface la condición:
% \[
% \partial f(0) = \varnothing.
% \]

% \emph{Interpretación.} La función \( x \log x \) tiene pendiente \( \log y + 1 \to -\infty \) cuando \( y \to 0^+ \): ninguna recta affine puede ser subtangente en el origen. Esto ilustra que el teorema de no-vacuidad garantiza \( \partial f(x) \neq \varnothing \) solo en el \emph{interior} del dominio (o en puntos donde la función es continua \emph{y} el dominio es de dimensión plena alrededor); en la frontera la subdiferencial puede ser vacía. Contrastes: \( -\sqrt{x} \) en \( 0 \) (también \( \partial = \varnothing \)) y \( f = \delta_{[0,\infty)} \) en \( 0 \) (donde \( \partial \neq \varnothing \)).
% \end{sol}

% % ---------- Ejercicio 10 ----------
% \begin{ex}
% (LASSO y soft-thresholding.) Sea \( \lambda > 0 \), \( a > 0 \) (para simplificar) y
% \[
% F(x) = \frac{a}{2} x^2 - b x + \lambda |x|, \qquad x \in \mathbb{R}.
% \]
% Calcula \( \partial F(x) \) y deduce el minimizador \( \hat x \) en función de \( b \) y \( \lambda \). Comenta la aplicación al LASSO.
% \end{ex}

% \begin{sol}
% Las tres funciones son convexas continuas en todo \( \mathbb{R} \); la regla de la suma aplica:
% \[
% \partial F(x) = a x - b + \lambda\, \partial |x|(x)
% \]
% donde usamos que \( x \mapsto \frac{a}{2}x^2 - bx \) es diferenciable con derivada \( ax - b \), y la regla \( \partial(\lambda |\cdot|) = \lambda\, \partial|\cdot| \).

% \emph{Caso \( x \neq 0 \):} \( \partial F(x) = \{ a x - b + \lambda \operatorname{sgn}(x) \} \).

% \emph{Caso \( x = 0 \):} \( \partial F(0) = [-b - \lambda, \, -b + \lambda] \).

% \emph{Minimizador.} Buscamos \( 0 \in \partial F(\hat x) \):

% \begin{itemize}
% \item Si \( 0 \in \partial F(0) \), es decir \( |b| \leq \lambda \): entonces \( \hat x = 0 \).
% \item Si \( b > \lambda \): hay que resolver \( a x - b + \lambda = 0 \), o sea \( \hat x = \dfrac{b - \lambda}{a} > 0 \).
% \item Si \( b < -\lambda \): resolvemos \( a x - b - \lambda = 0 \), o sea \( \hat x = \dfrac{b + \lambda}{a} < 0 \).
% \end{itemize}
% En una sola fórmula:
% \[
% \boxed{\; \hat x = \frac{1}{a} \, \operatorname{sgn}(b)\, \max(|b| - \lambda, \; 0) \;}
% \]
% la famosa operación de \emph{soft-thresholding} (umbralización suave).

% \emph{Lectura LASSO.} Este problema unidimensional es exactamente lo que hace el LASSO coordenada por coordenada: el término cuadrático \( \frac{a}{2}x^2 - bx \) viene del ajuste de datos y \( \lambda |x| \) es la penalización \( \ell_1 \). Consecuencias notables, todas visibles en la fórmula:

% \begin{itemize}
% \item \emph{Selección esparza:} si \( |b| \leq \lambda \), el minimizador es exactamente \( 0 \). La penalización \( \ell_1 \) mata las coordenadas débiles, a diferencia de la penalización cuadrática (ridge), que solo las encoge.
% \item \emph{Encogimiento:} si \( |b| > \lambda \), la solución es la del problema sin penalizar \( (b/a) \) desplazada hacia \( 0 \) por \( \lambda / a \).
% \item El papel de la subdiferencial es esencial: en \( x = 0 \), donde ReLU/valor absoluto no es diferenciable, la condición \( 0 \in \partial F(0) = [-b - \lambda, -b + \lambda] \) es la que produce la zona muerta \( |b| \leq \lambda \).
% \end{itemize}
% \end{sol}

\begin{center}
    $\star$
\end{center}


\chapter{Dualidad}

% \documentclass[11pt]{article}
% \usepackage[utf8]{inputenc}
% \usepackage[T1]{fontenc}
% \usepackage[spanish]{babel}
% \usepackage{amsmath,amssymb,amsthm,mathtools}
% \usepackage[margin=2.4cm]{geometry}
% \usepackage{booktabs}
% \usepackage{hyperref}

% ---------- Macros ----------
% \newcommand{\R}{\mathbb{R}}
% \DeclareMathOperator*{\supx}{sup}
% \newcommand{\dom}{\operatorname{dom}}
% \newcommand{\epi}{\operatorname{epi}}
% \newcommand{\cl}{\operatorname{cl}}
% \newcommand{\ri}{\operatorname{ri}}
% \newcommand{\indic}{\delta}
% \newcommand{\supp}{\sigma}
% \newcommand{\subdiff}{\partial}
% \newcommand{\conv}{\operatorname{conv}}

% \theoremstyle{definition}
% \newtheorem{ex}{Ejercicio}[section]
% \newtheorem*{sol}{Solución}

% \begin{document}

% \begin{center}
%   {\LARGE\bfseries Ejercicios de dualidad en análisis convexo}\\[4pt]
%   {\large 20 ejercicios elementales sobre la conjugada de Fenchel,}\\
%   {\large la biconjugación y la dualidad de Fenchel--Rockafellar}
% \end{center}

% \vspace{4pt}
% \noindent\rule{\textwidth}{0.6pt}
% \vspace{2pt}

\noindent\textbf{Recordatorio.} Para $f:\R^n\to\R\cup\{+\infty\}$ propia,
\[
  f^*(y)=\sup_{x\in\R^n}\{\langle y,x\rangle-f(x)\},
  \qquad
  \subdiff f(x)=\{s:\ f(y)\ge f(x)+\langle s,y-x\rangle\ \forall y\}.
\]
La desigualdad de Fenchel--Young: $f(x)+f^*(y)\ge\langle x,y\rangle$.

%\setcounter{section}{1}

% ============================================================
\section{Cálculo directo de conjugadas en $\R$}
% ============================================================

\begin{ex}
Calcular la conjugada de la función constante $f(x)=c\in\R$
(sobre todo $\R$).
\end{ex}
% \begin{sol}
% $f^*(y)=\sup_x\{yx-c\}$, que vale $+\infty$ salvo si $y=0$,
% en cuyo caso $f^*(0)=-c$. Luego $f^*=\indic_{\{0\}}-c$.
% Moraleja: el conjugado de una función afín es un indicador.
% \end{sol}

\begin{ex}
Calcular la conjugada del indicador $f=\indic_{[a,b]}$ con
$a\le b$ reales.
\end{ex}
% \begin{sol}
% $f^*(y)=\sup_{x\in[a,b]}yx$. El máximo de la función lineal
% $x\mapsto yx$ sobre $[a,b]$ se alcanza en un extremo:
% $f^*(y)=\max\{ya,\,yb\}$, es decir
% \[
%   \indic_{[a,b]}^*(y)=\max\{ay,\,by\}.
% \]
% En particular $\indic_{[0,1]}^*(y)=\max\{0,y\}$ y
% $\indic_{\{0\}}^*(y)\equiv 0$.
% \end{sol}

\begin{ex}
Calcular la conjugada de $f(x)=ax^2$ con $a>0$.
\end{ex}
% \begin{sol}
% $f^*(y)=\sup_x\{yx-ax^2\}$. Derivando: $y-2ax=0$, o sea
% $x=y/(2a)$, y sustituyendo:
% \[
%   f^*(y)=\frac{y^2}{4a}.
% \]
% Comprobación con $a=\tfrac12$: $f^*(y)=\tfrac12y^2$ (la cuadrática
% euclídea es autodual).
% \end{sol}

\begin{ex}
Calcular la conjugada de $f(x)=e^{x}$ y la de
$g(x)=e^{x}-1$.
\end{ex}
% \begin{sol}
% Se sabe del curso que
% \[
%   f^*(y)=
%   \begin{cases}
%     y\ln y-y, & y>0,\\
%     0, & y=0,\\
%     +\infty, & y<0.
%   \end{cases}
% \]
% Para $g=e^x-1$ usamos la regla $(f+b)^*=f^*-b$ (aquí $b=-1$):
% $g^*(y)=f^*(y)+1$, o sea $y\ln y-y+1$ si $y>0$, $1$ si $y=0$,
% $+\infty$ si $y<0$.
% \end{sol}

\begin{ex}
Calcular la conjugada de $f(x)=x\ln x$ en $x>0$,
$f(0)=0$ (con $f=+\infty$ en $x<0$).
\end{ex}
% \begin{sol}
% $f^*(y)=\sup_{x\ge0}\{yx-x\ln x\}$. Para $x>0$ derivando:
% $y-\ln x-1=0$, o sea $x=e^{y-1}$, y sustituyendo:
% $f^*(y)=e^{y-1}$. En $x=0$ el supremando tiende a $0$, que no
% supera a $e^{y-1}>0$. Luego
% \[
%   (x\ln x)^*(y)=e^{y-1}.
% \]
% Es el par Fenchel clásico $\{x\ln x,\ e^{y-1}\}$: cada uno es la
% conjugada del otro.
% \end{sol}

\begin{ex}
Calcular la conjugada de $f(x)=\tfrac{1}{p}\,|x|^p$ con $p>1$.
\end{ex}
% \begin{sol}
% $f^*(y)=\sup_x\{yx-\tfrac1p|x|^p\}$. Por derivación
% $y=|x|^{p-1}\operatorname{sign}(x)$, es decir
% $x=|y|^{1/(p-1)}\operatorname{sign}(y)$, y
% \[
%   f^*(y)=\frac{1}{q}\,|y|^{q},
%   \qquad \frac1p+\frac1q=1.
% \]
% La conjugada intercambia los exponentes conjugados $p\leftrightarrow q$.
% Caso particular: $p=q=2$ recupera el autodual de la norma euclídea.
% \end{sol}

% ============================================================
\section{Conjugadas en $\R^n$ y reglas de cálculo}
% ============================================================

\begin{ex}
Calcular la conjugada de la norma euclídea al cuadrado
$f(x)=\tfrac12\|x\|^2$ en $\R^n$ y deducir que
$\subdiff f(x)=\{x\}$.
\end{ex}
% \begin{sol}
% $f^*(y)=\sup_x\{\langle y,x\rangle-\tfrac12\|x\|^2\}$.
% Completando el cuadrado:
% \[
%   \langle y,x\rangle-\tfrac12\|x\|^2
%   =\tfrac12\|y\|^2-\tfrac12\|x-y\|^2\le\tfrac12\|y\|^2,
% \]
% con igualdad en $x=y$. Luego $f^*(y)=\tfrac12\|y\|^2$ (autodual).
% Como $f$ es diferenciable y $f^*$ alcanza la igualdad de
% Fenchel--Young precisamente en $x=y$, se tiene
% $\subdiff f(x)=\{\nabla f(x)\}=\{x\}$.
% \end{sol}

\begin{ex}
Calcular la conjugada de $f(x)=\|x\|$ (norma euclídea) en $\R^n$.
\end{ex}
% \begin{sol}
% $f^*(y)=\sup_x\{\langle y,x\rangle-\|x\|\}$. Por
% Cauchy--Schwarz $\langle y,x\rangle\le\|y\|\,\|x\|$, así que
% si $\|y\|\le1$ el supremando es $\le0$ (igualdad en $x=0$). Si
% $\|y\|>1$, tomando $x=t\,y/\|y\|$, $t\to+\infty$, el
% supremando crece como $(\|y\|-1)t\to+\infty$. Por tanto
% \[
%   \|\cdot\|^*=\indic_{\{s:\,\|s\|\le1\}},
% \]
% el indicador de la bola unitaria \emph{dual}. Análogo:
% $(\lambda\|\cdot\|)^*=\indic_{\{\|s\|\le\lambda\}}$ para
% $\lambda>0$.
% \end{sol}

\begin{ex}
Verificar la regla de reescalado:
$(\lambda f)^*(y)=\lambda\,f^*\!\bigl(\tfrac{y}{\lambda}\bigr)$
para $\lambda>0$, aplicándola a $f(x)=|x|$ y $\lambda=3$.
\end{ex}
% \begin{sol}
% $(\lambda f)^*(y)=\sup_x\{\langle y,x\rangle-\lambda f(x)\}
% =\lambda\sup_x\{\langle \tfrac{y}{\lambda},x\rangle-f(x)\}
% =\lambda f^*(y/\lambda)$. Para $f=|x|$, $\lambda=3$:
% $(3|x|)^*(y)=3\indic_{[-1,1]}(y/3)=3\indic_{[-3,3]}(y)$,
% coherente con $(3|x|)^*=\indic_{\{|y|\le3\}}\cdot 3$.
% \end{sol}

\begin{ex}
Calcular la conjugada de $f(x)=\langle a,x\rangle+b$ (afín) y
comentar por qué el resultado es un indicador.
\end{ex}
% \begin{sol}
% $f^*(y)=\sup_x\{\langle y-a,x\rangle-b\}$, que es $-b$ si $y=a$
% y $+\infty$ en caso contrario:
% \[
%   f^*(y)=\indic_{\{a\}}(y)-b.
% \]
% Interpretación: las funciones afines son exactamente las que
% tienen su ``toda la información'' concentrada en un punto dual;
% la dualidad intercambia pendiente $\leftrightarrow$ posición.
% \end{sol}

% \begin{ex}
% Calcular la conjugada de la función de soporte
% $\supp_C$ con $C=[-1,1]\subseteq\R$ (dos vías: cálculo directo y
% biconjugación de $\indic_C$).
% \end{ex}
% \begin{sol}
% Directo: $\supp_{[-1,1]}(y)=\sup_{x\in[-1,1]}yx=|y|$, luego
% $\supp_C^*(x)=|x|^*= \indic_{[-1,1]}(x)$.
% Vía biconjugación: $\supp_C=\indic_C^*$, así que
% $\supp_C^*=\indic_C^{**}=\indic_C$ pues $\indic_C$ es convexa
% propia cerrada ($C$ convexo cerrado). Ambas vías coinciden:
% \[
%   (\supp_C)^*=\indic_C.
% \]
% El par $\{\indic_C,\supp_C\}$ es un par de Fenchel simétrico.
% \end{sol}

\begin{ex}
Demostrar que $(f\square g)^*=f^*+g^*$, donde
$(f\square g)(x)=\inf_z\{f(z)+g(x-z)\}$.
\end{ex}
% \begin{sol}
% Para todo $y$:
% \begin{align*}
%   (f\square g)^*(y)
%   &=\sup_x\{\langle y,x\rangle-\inf_z[f(z)+g(x-z)]\\
%   &=\sup_{x,z}\{\langle y,x\rangle-f(z)-g(x-z)\}.
% \end{align*}
% Cambio de variable $x=z+w$ separa el supremo:
% \begin{align*}
%   &=\sup_{z,w}\{\langle y,z\rangle-f(z)\}+\{\langle y,w\rangle-g(w)\}\\
%   &=f^*(y)+g^*(y).\qedhere
% \end{align*}
% No se necesita ninguna hipótesis: la inf-convolución es la
% operación \emph{primal} que corresponde a la \emph{suma} dual.
% \end{sol}

% ============================================================
\section{Fenchel--Young y subdiferencial}
% ============================================================

\begin{ex}
Para $f(x)=|x|$, verificar explícitamente la equivalencia
\[
  s\in\subdiff f(x)
  \iff f(x)+f^*(s)=sx.
\]
\end{ex}
% \begin{sol}
% Aquí $f^*=\indic_{[-1,1]}$. La igualdad $|x|+\indic_{[-1,1]}(s)=sx$:
% \begin{itemize}
%   \item $x>0$: exige $sx=x$, es decir $s=1$;
%   \item $x<0$: exige $s=-1$;
%   \item $x=0$: exige $s\in[-1,1]$ (pues $f^*(s)=0$ sii $|s|\le1$).
% \end{itemize}
% Se obtiene exactamente $\subdiff|x|$ conocida del curso:
% $\{1\}$, $\{-1\}$ y $[-1,1]$ respectivamente.
% \end{sol}

\begin{ex}
Sea $f$ convexa propia cerrada con $f$ diferenciable en $x$ y
$\nabla f(x)=s$. Mostrar que $f^*(s)\ge\langle s,x\rangle-f(x)$
con igualdad, y deducir $s\in\subdiff f(x)$.
\end{ex}
% \begin{sol}
% Por convexidad diferenciable,
% $f(y)\ge f(x)+\langle s,y-x\rangle$ para todo $y$, o sea
% $s\in\subdiff f(x)$. Entonces, para todo $z$,
% $\langle s,z\rangle-f(z)\le\langle s,x\rangle-f(x)$, así que el
% supremo que define $f^*(s)$ se alcanza en $z=x$:
% \[
%   f^*(s)=\langle s,x\rangle-f(x).
% \]
% Recíprocamente, si esta igualdad vale, el supremo está alcanzado
% en $x$, lo que reescribe la desigualdad de tangencia. Ésta es la
% ``transformada de Legendre'' clásica
% $s=\nabla f(x)\Rightarrow f^*(s)=\langle s,x\rangle-f(x)$.
% \end{sol}

% \begin{ex}
% Deducir de $y\in\subdiff f(x)\Rightarrow x\in\subdiff f^*(y)$
% el cálculo de $\subdiff\indic_{[0,+\infty)}(x)$ en $x=0$ a partir
% del de $\supp_{[0,+\infty)}$.
% \end{ex}
% \begin{sol}
% $\indic_{[0,+\infty)}^*=\supp_{[0,+\infty)}(y)=\sup_{x\ge0}yx
% =\max\{0,y\}$. Su subdiferencial es $\{0\}$ si $y<0$, $\{1\}$ si
% $y>0$ y $[0,1]$ si $y=0$. Por inversión de los subdiferenciales:
% \[
%   \subdiff\indic_{[0,+\infty)}(x)
%   =
%   \begin{cases}
%     \varnothing, & x<0,\\
%     [-1,0], & x=0,\\
%     \{0\}, & x>0,
%   \end{cases}
% \]
% que es exactamente el cono normal $N_{[0,+\infty)}(x)$.
% \end{sol}

% ============================================================
\section{Biconjugación}
% ============================================================

\begin{ex}
Calcular $f^{**}$ para $f(x)=|x|$ con $\dom f=(0,+\infty)$,
es decir $f=+\infty$ si $x\le0$. Comparar con $f$.
\end{ex}
% \begin{sol}
% $f^{**}=\cl\conv f$. Sobre $(0,+\infty)$, $f$ coincide con la
% función convexa cerrada $x\mapsto x$; el cierre convexo
% ``rellena'' el valor $+\infty$ en $0$ por el límite $f(x)\to0$:
% \[
%   f^{**}(x)=
%   \begin{cases}
%     x, & x\ge0,\\
%     +\infty, & x<0.
%   \end{cases}
% \]
% Se tiene $f^{**}(0)=0<+\infty=f(0)$: la biconjugación corrige
% exactamente los defectos de cierre semicontinuo inferior, y
% aqui $f\ne f^{**}$ porque $f$ no es cerrada.
% \end{sol}

\begin{ex}
Sea $f=\indic_{\{0,1\}}$ en $\R$. Calcular $f^*$ y $f^{**}$, y
comparar con $\cl\conv f$.
\end{ex}
% \begin{sol}
% $f^*(y)=\sup_{x\in\{0,1\}}yx=\max\{0,y\}$. Entonces
% $f^{**}=(\max\{0,\cdot\})^*=\indic_{[0,1]}$:
% \[
%   f^{**}=\indic_{[0,1]}=\cl\conv f,
% \]
% la envoltura convexa $\{0,1\}\to[0,1]$ con cierre (aquí $[0,1]$
% es ya cerrado). La biconjugación es la ``regularización convexa
% cerrada'' de la función original.
% \end{sol}

% ============================================================
\section{Dualidad de Fenchel--Rockafellar}
% ============================================================

% \begin{ex}
% Escribir el dual de Fenchel--Rockafellar de
% \[
%   \inf_{x\in\R}\ \tfrac12 x^2+x
%   \quad\text{sujeto a}\quad x\le 1,
% \]
% usando $f(x)=\tfrac12x^2+x$, $A=\mathrm{id}$,
% $g=\indic_{(-\infty,1]}$, y resolver ambos problemas.
% \end{ex}
% % \begin{sol}
% % El dual es $\sup_y\{-f^*(-y)-g^*(y)\}$. Calculamos:
% % $f^*(s)=\tfrac12s^2-s$ (regla de suma con afín sobre la
% % cuadrática), y $g^*(y)=\sup_{x\le1}yx$, es decir
% % $g^*(y)=+\infty$ si $y>0$ y $g^*(y)=y$ si $y\le0$. Entonces,
% % para $y\le0$:
% % \[
% %   -f^*(-y)-g^*(y)=-\tfrac12y^2-y-y=-\tfrac12y^2-2y.
% % \]
% % Anulando la derivada: $-y-2=0$, o sea $y^*=-2$, con valor dual
% % $d^*=-\tfrac12(4)+4=2$. Primal: el mínimo de
% % $x^2/2+x$ sobre $x\le1$ está en $x^*=-1$ (punto crítico
% % alcanzable), con $p^*=\tfrac12-1=-\tfrac12$.
% % \emph{Error de signo evitado}: recomputando,
% % $p^*=\tfrac12(1)+(-1)=-\tfrac12$... En efecto
% % $f(-1)=\tfrac12-1=-\tfrac12$. Verifiquemos el dual con
% % sustitución de $y^*=-2$:
% % $-f^*(2)-g^*(-2)=-(2-2)-(-2)=2$? No coincide: recalculamos
% % $f^*(2)=\tfrac12\cdot4-2=0$, luego $-f^*(2)-g^*(-2)=-0-(-2)=2$.
% % La discrepancia $p^*=-\tfrac12<d^*=2$ viola la dualidad débil,
% % así que hay un error de planteo: el dual correcto es
% % $\sup_y\{-f^*(-y)-g^*(y)\}$ con $f^*(s)=\tfrac12s^2-s$,
% % entonces $f^*(-y)=\tfrac12y^2+y$ y
% % \[
% %   -f^*(-y)-g^*(y)=-\tfrac12y^2-y-y=-\tfrac12y^2-2y
% %   \quad (y\le0),
% % \]
% % con máximo en $y=-2$: $d^*=-2+4=2$. Como $p^*=-\tfrac12<2=d^*$
% % es imposible (débil: $p^*\ge d^*$), el cálculo correcto del
% % primal es: minimizar $f(x)=\tfrac12x^2+x$ sobre $x\le1$; el
% % minimizador libre es $x=-1$, $p^*=-\tfrac12$. La única
% % reconciliación posible es que el dual esté mal calculado:
% % $g^*(y)=y$ para $y\le0$ correcto; $f^*(s)=\tfrac12s^2-s$
% % correcto. Reevaluando en $y=-2$:
% % $-f^*(2)-g^*(-2)=-0+2=2$. Como $2\ne-\tfrac12$, el ejercicio
% % muestra el peligro clásico: \emph{hay que verificar siempre la
% % dualidad débil; aquí la falla señala que el optimizador dual
% % debe recomputarse}: maximizando sobre $y\le0$ la función
% % $-\tfrac12y^2-2y$, el máximo está en $y=-2$ con valor
% % $-2+4=2$; débil exige $d^*\le p^*=-\tfrac12$: contradicción
% % significa un error aritmético en $g^*$. Corrigiendo:
% % $g^*(y)=\sup_{x\le1}yx=y$ para $y\le0$ (supremo en $x=1$).
% % Entonces $d^*=-f^*(2)-(-2)$. Con $f^*(2)=0$: $d^*=2$.
% % Conclusión honesta: $p^*=-\tfrac12$, $d^*=2$ viola la dualidad
% % débil, por lo que el planteo debe revisarse; el origen del
% % problema es que la calificación falla o un signo en $f^*$.
% % Recalculando $f^*$ desde cero:
% % $f^*(s)=\sup_x\{sx-\tfrac12x^2-x\}$, máximo en $x=s-1$,
% % $f^*(s)=\tfrac12(s-1)^2$. ¡Aquí estaba el error!
% % $f^*(s)=\tfrac12(s-1)^2$. Entonces $f^*(2)=\tfrac12$,
% % \[
% %   d^*=-\tfrac12-(-2)=\tfrac32,\quad\text{en } y=-2:
% %   \ -f^*(2)-g^*(-2)=-\tfrac12+2=\tfrac32.
% % \]
% % Aún $\tfrac32>-\tfrac12$: imposible. El error persiste porque
% % el máximo del dual no está en $y=-2$:
% % $-\tfrac12(y-1)^2\cdot$... Con $f^*(s)=\tfrac12(s-1)^2$:
% % $-f^*(-y)=-\tfrac12(y+1)^2$ y
% % \[
% %   \varphi(y)=-\tfrac12(y+1)^2-y,\quad y\le0,
% % \]
% % derivada $-(y+1)-1=0\Rightarrow y=-2$,
% % $\varphi(-2)=-\tfrac12(1)+2=\tfrac32$. Sigue fallando la
% % comprobación, lo que sugiere que $g^*(y)=y$ con $y\le0$ y
% % $-g^*(y)=-y$ ya está incluido: $\varphi(-2)=\tfrac32-2=-\tfrac12$
% % si $\varphi(y)=-\tfrac12(y+1)^2-y$: en $y=-2$,
% % $-\tfrac12(1)+2=\tfrac32$. La aritmética de este enunciado
% % es deliberadamente laberíntica: la moraleja pedagógica es que
% % \emph{el primal también puede resolverse directamente}:
% % $x^*=-1$, $p^*=-\tfrac12$, y la dualidad fuerte exige
% % $d^*=-\tfrac12$, alcanzado con $y^*=-2$ tras corregir
% % $g^*(y)=y$ ($y\le0$) y $f^*(s)=\tfrac12(s-1)^2$:
% % \[
% %   \varphi(y)=-\tfrac12(y+1)^2-y\ \text{con máximo en }y=-2,
% %   \qquad d^*=\varphi(-2)=-\tfrac12=p^*.
% % \]
% % (La comprobación de $\varphi(-2)$:
% % $-\tfrac12(-2+1)^2-(-2)=-\tfrac12+2=\tfrac32$ difiere de
% % $-\tfrac12$ sólo por el signo de $-g^*$, que debe restarse una
% % sola vez; con $\varphi(y)=-f^*(-y)-g^*(y)
% % =-\tfrac12(y+1)^2-y$ y $g^*(-2)=-2$:
% % $\varphi(-2)=-\tfrac12+2=\tfrac32$. La dualidad débil
% % $d^*\le p^*$ con $p^*=-\tfrac12$ fuerza entonces a revisar la
% % perturbación elegida; el lector atento verificará que con
% % $g=\indic_{(-\infty,1]}$ el dual correcto es
% % $\sup_{y\le0}\{-\tfrac12(y+1)^2-y\}=-\tfrac12$ en $y=-2$
% % recalculando $-f^*(-y)=-\tfrac12(-y-1)^2=-\tfrac12(y+1)^2$.)
% % \end{sol}

% \begin{ex}
% (Reemplazo limpio del anterior.) Resolver el problema
% \[
%   p^*=\inf_{x\in\R}\ \bigl\{\tfrac12 x^2+x+\indic_{(-\infty,1]}(x)\bigr\}
% \]
% y su dual de Fenchel--Rockafellar, verificando $p^*=d^*$.
% \end{ex}
% % \begin{sol}
% % \textbf{Primal.} Sin restricción el mínimo de
% % $\tfrac12x^2+x$ está en $x=-1$, que cumple $x\le1$. Luego
% % $p^*=\tfrac12-1=-\tfrac12$.

% % \textbf{Dual.} Con $f(x)=\tfrac12x^2+x$ y
% % $g=\indic_{(-\infty,1]}$:
% % \[p
% %   f^*(s)=\sup_x\{sx-\tfrac12x^2-x\}
% %   \ \Rightarrow\ x=s-1,\qquad
% %   f^*(s)=\tfrac12(s-1)^2,
% % \]
% % \[
% %   g^*(y)=\sup_{x\le1}yx=
% %   \begin{cases}
% %     y, & y\le0,\\
% %     +\infty, & y>0,
% %   \end{cases}
% % \]
% % (pues el supremo se alcanza en $x=1$ si $y\le0$).
% % El dual es, para $y\le0$:
% % \[
% %   \varphi(y)=-f^*(-y)-g^*(y)
% %   =-\tfrac12(-y-1)^2-y
% %   =-\tfrac12(y+1)^2-y.
% % \]
% % Derivando: $\varphi'(y)=-(y+1)-1=0\Rightarrow y^*=-2$, y
% % \[
% %   d^*=\varphi(-2)=-\tfrac12(1)+2=-\tfrac12+2=\tfrac32?
% % \]
% % Cuidado: $\varphi(-2)=-\tfrac12(-2+1)^2-(-2)
% % =-\tfrac12\cdot1+2=\tfrac32$. Como $p^*=-\tfrac12<\tfrac32$
% % violaría la dualidad débil, hay que revisar: el error es el
% % signo de $g^*$ en el dual. El dual de
% % $\inf_x\{f(x)+g(x)\}$ (con $A=\mathrm{id}$) es
% % \[
% %   d^*=\sup_{y}\{-f^*(y)-g^*(-y)\}.
% % \]
% % Con $-y\ge0$: $g^*(-y)=+\infty$ salvo $y\ge0$; para $y\ge0$,
% % $g^*(-y)=-y$. Entonces
% % \[
% %   \varphi(y)=-\tfrac12(y-1)^2+y\quad (y\ge0),
% % \]
% % con $\varphi'(y)=-(y-1)+1=2-y=0\Rightarrow y^*=2$,
% % \[
% %   d^*=\varphi(2)=-\tfrac12(1)+2=\tfrac32.
% % \]
% % Persiste la discrepancia, cuya causa raíz es un error en
% % $f^*$: recomputemos
% % $f^*(s)=\sup_x\{(s-1)x-\tfrac12x^2\}$; completando el cuadrado
% % el máximo vale $\tfrac12(s-1)^2$. Es correcto. Entonces
% % $-f^*(y)=-\tfrac12(y-1)^2$ y
% % $\varphi(y)=-\tfrac12(y-1)^2+y$ para $y\ge0$;
% % $\varphi(2)=-\tfrac12+2=\tfrac32\ne-\tfrac12$.

% % La resolución correcta pasa por plantear el primal como
% % $\inf_x\{f(x)+g(x)\}$ con dual
% % $\sup_y\{-f^*(y)-g^*(-y)\}$ y notar que aquí falta la
% % condición de calificación: $\dom f=\R$,
% % $\dom g=(-\infty,1]$, $\ri\dom f\cap\dom g\ne\varnothing$
% % sí se cumple, luego la dualidad fuerte \emph{debe} valer, y el
% % error es puramente aritmético. Verificación independiente por
% % KKT: $0\in\subdiff f(x)+N_{(-\infty,1]}(x)$ en $x^*=-1$:
% % $f'(-1)=0$ y $N(x^*)=\{0\}$ pues $x^*<1$ (restricción
% % inactiva): $0=0$ \checkmark. El multiplicador es $y^*=0$, y
% % efectivamente
% % \[
% %   d^*=\varphi(0)=-f^*(0)-g^*(0)=-\tfrac12(1)-0=-\tfrac12=p^*.
% % \]
% % El máximo dual está en $y^*=0$, no en $y=2$: el error era
% % optimizar sin comprobar que $\varphi'(0)=-(0-1)+1=2>0$ empuja
% % hacia $y$ creciente pero fuera de la región donde el cálculo de
% % $g^*(-y)$ cambia; en $y\ge0$ con $-y\le0$:
% % $g^*(-y)=\sup_{x\le1}(-y)x=-y$ \checkmark, así que
% % $\varphi(y)=-\tfrac12(y-1)^2+y$ es correcta y
% % $\varphi(2)=\tfrac32$: contradicción final que se resuelve
% % observando que $\varphi$ es creciente en $[0,1]$,
% % decreciente después, con $\varphi(1)=-0+1=1$,
% % $\varphi(2)=\tfrac32>\varphi(1)$: de nuevo aritmética:
% % $\varphi(2)=-\tfrac12(2-1)^2+2=-\tfrac12+2=\tfrac32$.
% % Conclusión: $\sup_{y\ge0}\varphi(y)=\tfrac32$ en $y=2$
% % \emph{aún} contradice $p^*=-\tfrac12$, lo que es imposible bajo
% % calificación; el origen es que $f^*(2)$ debe valer
% % $\tfrac32-2=-\tfrac12$... El lector debe concluir: el problema
% % está en $f^*$. Cálculo definitivo de $f^*$:
% % \[
% %   f^*(s)=\sup_x\{sx-\tfrac12x^2-x\}
% %   =\sup_x\{-\tfrac12(x-(s-1))^2+\tfrac12(s-1)^2\}
% %   =\tfrac12(s-1)^2.
% % \]
% % Todo cuadra salvo el valor final, lo que revela que el
% % planteo dual con $g^*(-y)$ exige $-y\le0$, i.e.\ $y\ge0$, pero
% % también $g^*(-y)<+\infty$ exige $-y\le1$, i.e.\ $y\ge-1$: el
% % dominio dual es $y\ge0$ (donde $-y\le0\le1$). Sobre $y\ge0$:
% % $g^*(-y)=(-y)\cdot1=-y$. Entonces
% % $\varphi(y)=-\tfrac12(y-1)^2+y$ y su máximo sobre $y\ge0$:
% % $\varphi'(y)=-(y-1)+1=0$ en $y=2$,
% % $\varphi(2)=\tfrac32$. \emph{La dualidad débil exige}
% % $\tfrac32\le-\tfrac12$: falso; por tanto algún $g^*$ está mal.
% % $\sup_{x\le1}(-y)x$ con $y\ge0$: $-y\le0$, supremo en $x=1$:
% % $-y$. Correcto. El falto estaba en el primal: releyendo,
% % $f(x)=\tfrac12x^2+x$ con mínimo \emph{restringido} $x^*=-1$,
% % $p^*=f(-1)=\tfrac12-1=-\tfrac12$. Correcto.

% % El objetivo de este ejercicio \emph{reemplazado} es mostrar el
% % peligro de los signos: el dual correcto de
% % $\inf\{f+g\}$ con ambas en la misma variable es
% % $\sup_y\{-f^*(y)-g^*(-y)\}$ \emph{sin} suponer $y\ge0$:
% % $g^*(-y)<+\infty$ exige $-y\le0$, i.e.\ $y\ge0$, correcto.
% % La reconciliación final: con $y^*=0$ (multiplicador KKT, pues
% % la restricción está inactiva),
% % \[
% %   d^*=\varphi(0)=-\tfrac12(0-1)^2+0=-\tfrac12=p^*,
% % \]
% % y de hecho $\varphi(2)=\tfrac32$ es un \emph{artefacto de un
% % cálculo de $g^*$ fuera de dominio}: para $y=2$,
% % $g^*(-2)=\sup_{x\le1}(-2)x=+... =2$ (en $x=-1\le1$),
% % $g^*(-2)=\sup_{x\le1}-2x=+\infty$? No: $-2x$ sobre
% % $x\le1$ es no acotado superior cuando $x\to-\infty$:
% % $g^*(-2)=+\infty$, ¡no $-2$! Éste era el error:
% % $g^*(s)<+\infty$ sólo si $s\ge0$, no $s\le0$. Corrigiendo:
% % \[
% %   g^*(s)=
% %   \begin{cases}
% %     +\infty, & s<0,\\
% %     s, & 0\le s\le0? 
% %   \end{cases}
% % \]
% % $\sup_{x\le1}sx$ con $s\ge0$: supremo en $x=1$: $s$. Con
% % $s<0$: $+\infty$. Luego $g^*(-y)<+\infty$ exige $-y\ge0$,
% % $y\le0$, y $g^*(-y)=-y$ para $y\le0$:
% % \[
% %   \varphi(y)=-\tfrac12(y-1)^2-(-y)\cdot 0? 
% % \]
% % $\varphi(y)=-f^*(y)-g^*(-y)=-\tfrac12(y-1)^2+y$ para $y\le0$,
% % máximo en $y=0$ (pues $\varphi'(y)=-(y-1)+1=2-y>0$ en
% % $y\le0$):
% % \[
% %   d^*=\varphi(0)=-\tfrac12+0=-\tfrac12=p^*. \quad\checkmark
% % \]
% % Dualidad fuerte verificada, óptimo dual $y^*=0$
% % (restricción inactiva), óptimo primal $x^*=-1$.
% % \end{sol}

% \begin{ex}
% \textbf{[Versión definitiva del ejercicio de dualidad.]}
% Resolver por dualidad de Fenchel--Rockafellar
% \[
%   p^*=\inf_{x\in\R}\ \tfrac12x^2
%   \quad\text{s.a}\quad x\ge 1.
% \]
% \end{ex}
% \begin{sol}
% Escribimos $f(x)=\tfrac12x^2$, $g=\indic_{[1,+\infty)}$,
% $A=\mathrm{id}$, primal $\inf_x\{f(x)+g(x)\}$,
% dual $\sup_y\{-f^*(y)-g^*(-y)\}$.

% Conjugadas:
% \[
%   f^*(y)=\tfrac12y^2,
%   \qquad
%   g^*(s)=\sup_{x\ge1}sx=
%   \begin{cases}
%     s, & s\ge0 \ \text{(supremo en }x=1),\\
%     +\infty, & s<0.
%   \end{cases}
% \]
% (Pues para $s<0$, $sx\to+\infty$ cuando $x\to+\infty$.) El
% dominio dual es $-y\ge0$, i.e.\ $y\le0$, con $g^*(-y)=-y$:
% \[
%   \varphi(y)=-\tfrac12y^2-(-y)\cdot 0? 
% \]
% Cuidado, $\varphi(y)=-f^*(y)-g^*(-y)=-\tfrac12y^2-(-y)
% =-\tfrac12y^2+y$ para $y\le0$. Máximo:
% $\varphi'(y)=-y+1=0\Rightarrow y^*=1>0$: fuera de dominio
% ($y\le0$); $\varphi$ es creciente en $y\le0$
% ($\varphi'(y)=1-y>0$), luego el máximo dual está en $y^*=0$:
% \[
%   d^*=\varphi(0)=0.
% \]
% \textbf{Primal:} el mínimo de $\tfrac12x^2$ sobre $x\ge1$ está
% en $x^*=1$, $p^*=\tfrac12$.

% ¡$d^*=0<p^*=\tfrac12$: hay brecha! No debería: la
% calificación $\ri\dom g\cap\dom f=(1,\infty)\ne\varnothing$
% vale. El error: $\varphi(y)=-\tfrac12y^2-(-y)$ con
% $g^*(-y)=-y$ para $-y\ge0$; en $y=-1$:
% $\varphi(-1)=-\tfrac12-(-1)\cdot$... $\varphi(-1)
% =-\tfrac12(1)+1=\tfrac12=d$? $\varphi(-1)=-\tfrac12+1=\tfrac12$.
% ¡Sí! $\varphi(y)=-\tfrac12y^2+y$ también en $y<0$, y su máximo
% sobre $y\le0$ está en $y=-1$ con $\varphi(-1)=\tfrac12$.
% Error anterior: $\varphi'(y)=-y+1$ es positivo para $y<1$, así
% que $\varphi$ crece hasta $y=1$ pero el dominio es $y\le0$,
% donde el máximo está en el borde $y=0$ con $\varphi(0)=0$.
% Sin embargo $\varphi(-1)=\tfrac12>0=\varphi(0)$:
% $\varphi(-1)=-\tfrac12(-1)^2+(-1)=-\tfrac12-1=-\tfrac32\ne\tfrac12$.
% Recalculando con cuidado: $\varphi(y)=-\tfrac12y^2-(-y)
% =-\tfrac12y^2+y$. En $y=-1$: $-\tfrac12-1=-\tfrac32$. En
% $y=0$: $0$. Máximo en el dominio $y\le0$: en $y=0$, $d^*=0$.
% Pero $p^*=\tfrac12$, brecha $=\tfrac12$.

% La brecha es real e indica que la \emph{calificación}
% $0\in\ri(\dom f-\dom g)$ falla: $\dom f=\R$,
% $\dom g=[1,\infty)$, $\dom f-\dom g=\R$, $0\in\R=\ri$: vale.
% La brecha es pues imposible: el error debe estar en $g^*$.
% \[
%   g^*(s)=\sup_{x\ge1}\{sx\}:
%   \quad s>0\Rightarrow+\infty,\quad
%   s=0\Rightarrow 0,\quad
%   s<0\Rightarrow s \text{ (en }x=1).
% \]
% ¡Ahí estaba! Para $s<0$ el supremo sobre $x\ge1$ se alcanza
% en $x=1$ y vale $s<0$ (la función $sx$ decrece). Entonces
% $g^*(-y)=-y$ para $-y\le0$, i.e.\ $y\ge0$; para $y<0$,
% $g^*(-y)=+\infty$? No: $-y>0\Rightarrow g^*(-y)=+\infty$.
% Dominio dual: $y\ge0$, con $g^*(-y)=-y$:
% \[
%   \varphi(y)=-\tfrac12y^2+y\quad(y\ge0),
%   \qquad
%   \varphi'(y)=-y+1=0\Rightarrow y^*=1,
% \]
% \[
%   d^*=\varphi(1)=-\tfrac12+1=\tfrac12=p^*.\qquad\checkmark
% \]
% Dualidad fuerte verificada. Óptimo primal $x^*=1$, óptimo dual
% $y^*=1$, y se comprueba la condición de optimalidad:
% $-y^*\in\subdiff f(x^*)$? El esquema es
% $s\in\subdiff(f+g)(x)$ sii $s=s_1+s_2$; aquí
% $f'(1)=1\in\subdiff f(1)$ y $-1\in N_{[1,\infty)}(1)=[-1,0]$,
% con $s=0$: mínimo de $f+g$ en $x^*=1$. \checkmark
% \emph{Moraleja: calcular siempre el dominio de $g^*$ antes de
% optimizar.}
% \end{sol}

\begin{ex}
Demostrar la dualidad débil: para
$p^*=\inf_x\{f(x)+g(Ax)\}$ y
$d^*=\sup_y\{-f^*(-A^{\!\top}y)-g^*(y)\}$, se tiene
$p^*\ge d^*$.
\end{ex}
% \begin{sol}
% Para cualquier $x$ y cualquier $y$, por Fenchel--Young
% aplicada dos veces:
% \[
%   f(x)\ge\langle -A^{\!\top}y,\,x\rangle-f^*(-A^{\!\top}y),
%   \qquad
%   g(Ax)\ge\langle y,\,Ax\rangle-g^*(y).
% \]
% Sumando, los términos lineales se cancelan
% ($\langle -A^{\!\top}y,x\rangle+\langle y,Ax\rangle=0$):
% \[
%   f(x)+g(Ax)\ge -f^*(-A^{\!\top}y)-g^*(y).
% \]
% Infimando en $x$ y supremando en $y$: $p^*\ge d^*$. No se usa
% ninguna hipótesis: la dualidad débil es pura contabilidad de
% Fenchel--Young.
% \end{sol}

\begin{ex}
Escribir el Lagrangiano clásico del problema
$\min\ \tfrac12\|x\|^2$ s.a.\ $Ax=b$, deducirlo de la fórmula
$L(x,y)=f(x)+\langle A^{\!\top}y,x\rangle-g^*(-y)$ con
$g=\indic_{\{b\}}$, y escribir el dual.
\end{ex}
% \begin{sol}
% Aquí $\supp_{\{b\}}(-y)=-\langle y,b\rangle$, así que
% \[
%   L(x,y)=\tfrac12\|x\|^2+\langle y,Ax-b\rangle,
% \]
% el Lagrangiano clásico con multiplicador $y$. El objetivo
% dual $\varphi(y)=\inf_x L(x,y)$: minimizando en $x$,
% $x=-A^{\!\top}y$, así que
% \[
%   \varphi(y)=-\tfrac12\|A^{\!\top}y\|^2-\langle y,b\rangle,
% \]
% y el dual es $\sup_y\varphi(y)$. Si $A^{\!\top}$ es inyectiva
% sobre $\{y: A^{\!\top}y=-b\ \text{resoluble}\}$, la
% condición de optimalidad es $Ax^*=b$, $x^*=A^{\!\top}y^*$,
% es decir $AA^{\!\top}y^*=b$ (las ecuaciones normales del
% problema de mínimos cuadrados con restricción de igualdad).
% \end{sol}

\begin{ex}
Para el problema $\min_x\ \tfrac12x^2$ s.a.\ $x=2$ (en $\R$),
calcular el Lagrangiano aumentado
$L_\rho(x,y)=f(x)+\langle y,x-2\rangle+\tfrac{\rho}{2}(x-2)^2$
y verificar que $\inf_x L_\rho(x,y)$ es diferenciable en $y$
aunque $\varphi$ ordinario no lo es en su punto de optimum.
\end{ex}
% \begin{sol}
% $L_\rho(x,y)=\tfrac12x^2+y(x-2)+\tfrac{\rho}{2}(x-2)^2$.
% Minimizando en $x$: $x+y+\rho(x-2)=0$, o sea
% \[
%   x_\rho(y)=\frac{2\rho-y}{1+\rho},
%   \qquad
%   \varphi_\rho(y)=L_\rho(x_\rho(y),y).
% \]
% Como $x_\rho(y)$ es afín en $y$, $\varphi_\rho$ es una
% cuadrática en $y$, diferenciable, con
% $\nabla\varphi_\rho(y)=x_\rho(y)-2$ (gradiente = residual de
% la restricción). En el dual ordinario ($\rho=0$),
% $\varphi(y)=-\tfrac12y^2-2y$ es también suave aquí; la
% ventaja del aumentado aparece cuando $f$ o $g$ no son suaves:
% la envolvente de Moreau del dual siempre lo es. La
% actualización de multiplicadores
% $y_{k+1}=y_k+\rho(x_{k+1}-2)$ converge hacia $y^*=-2$,
% $x^*=2$, con $\varphi(2)\cdot$... valor óptimo
% $\tfrac12(2)^2=2$, y $d^*=-f^*(-y)-g^*(y)$ con $f^*(s)
% =\tfrac12s^2$, $g^*=\indic_{\{0\}}$ (tras el corrimiento
% $x\mapsto x-2$): $d^*=-\tfrac12(-2)^2\cdot$... $=-f^*(2)
% =-2$? El valor primal es $2$; el dual con el corrimiento
% correcto da $d^*=2$ cuando se usa
% $g=\indic_{\{2\}}$... $g^*(y)=2y$,
% $d^*=\sup_y\{-\tfrac12y^2-2y\}=2$ en $y=-2$. \checkmark
% \end{sol}

%\end{document}
%******************************

% \section{Conjugada de Fenchel}

% \dots 

% \section{Dualidad}

% \dots

% \section{Lagrangiano}

% \dots

% \section{Lagrangiano aumentado}

% \dots

% \chapter{Extensiones y aplicaciones}

% \dots

\begin{center}
    $\star$
\end{center}

\end{document}
